\فصل{مروری بر ادبیات}
\label{ch:literature-review}

\قسمت{مقدمه و دامنهٔ مرور}
\label{sec:lit-review-introduction}

پیش‌بینی حرکت ذرات، به شناخت جریان و نیروهای مؤثر وابسته است.
همین حرکت، می‌تواند داده‌ای برای تخمین کمیت‌های پنهان فراهم کند.
بنابراین، حل مستقیم و بازسازی معکوس، دو جهت مرتبط پژوهش هستند.
در ریزجریان‌ها، انتخاب مدل به مقیاس‌ها و سازوکارهای غالب وابسته است.
کوچکی هندسه، به‌تنهایی همهٔ اثرهای اینرسی یا دیواره را حذف نمی‌کند.
\cite{SquiresQuake2005}

این فصل، پیشینهٔ پژوهش را حول چهار پرسش سازمان می‌دهد.
پرسش نخست، طراحی رهاسازی برای تشکیل آرایش چندذره‌ای است.
پرسش دوم، یادگیری پاسخ برای ورودی‌های تابعی متفاوت است.
پرسش سوم، بازسازی جریان از حرکت تعداد محدودی ذره است.
پرسش چهارم، طراحی تحریک مغناطیسی برای بازتولید مسیر هدف است.
هدف مرور، روشن‌کردن ارتباط و تفاوت این پرسش‌ها است.

شبکه‌های فیزیک‌آگاه، معادلات و داده‌ها را در تقریب پاسخ ترکیب می‌کنند.
یادگیری عملگر، تقریب نگاشت میان خانواده‌ای از ورودی‌ها و پاسخ‌ها را دنبال می‌کند.
در نتیجه، تقریب یک میدان و یادگیری خانوادهٔ میدان‌ها متفاوت هستند.
همچنین، یادگیری عملگر مستقیم با آموزش یک نگاشت معکوس یکسان نیست.
این تمایزها، مبنای بررسی کاربردهای مرتبط قرار می‌گیرند.
\cite{Raissi2019PINN,Kovachki2023NeuralOperator}

مقایسهٔ منابع، براساس مدل فیزیکی و اطلاعات مسئله انجام می‌شود.
نوع مشاهده، مجهول، پیش‌دانسته‌ها و مرجع ارزیابی نیز ثبت می‌شوند.
شمار ذرات، بدون مدل حرکت و اطلاعات زمانی، معیار کافی نیست.
همچنین، خطای مسیر، خطای میدان و خطای ورودی جدا تفسیر می‌شوند.
مطالعات بنیادی، برای تعریف مفاهیم و فرض‌های مدل حفظ شده‌اند.
مطالعات جدیدتر، برای بررسی توسعه‌ها و محدودیت‌های روش‌ها به‌کار می‌روند.
این مرور، ادعای پوشش نظام‌مند تمام پایگاه‌های پژوهشی ندارد.

ابتدا، مبانی ضروری جریان و حرکت ذره معرفی می‌شوند.
سپس، اصول شبکه‌های فیزیک‌آگاه و یادگیری عملگر بررسی می‌شوند.
در ادامه، طراحی رهاسازی و تشکیل آرایش مطرح می‌شود.
پس از آن، بازسازی جریان و تحریک از حرکت ذرات بررسی می‌شود.
فصل با تلفیق پیشینه و تعیین جایگاه پژوهش حاضر پایان می‌یابد.
جزئیات مدل‌های اجراشده و نتایج عددی، در فصل‌های بعد ارائه می‌شوند.

\قسمت{مبانی ضروری جریان و حرکت ذرات}
\label{sec:lit-flow-particle-foundations}

پیش‌بینی حرکت ذرات، به شناخت جریان و نیروهای واردشده نیاز دارد.
نوع مدل ذره نیز بر تفسیر مسیر مشاهده‌شده اثر می‌گذارد.
در این بخش، مبانی لازم برای بررسی پژوهش‌های مرتبط مرور می‌شوند.
تمرکز بر جریان کانالی، حرکت اینرسی‌دار و انتقال مغناطیسی است.

\زیرقسمت{جریان کانالی و مدل‌سازی جریان لزج}
\label{sec:lit-channel-flow}

اسکوایرز و کوئیک، فیزیک ریزجریان‌ها را با گروه‌های بی‌بعد بررسی کردند.
این چارچوب، اهمیت نسبی سازوکارهای فیزیکی را مشخص می‌کند.
در نتیجه، انتخاب مدل به مقیاس‌های مسئله وابسته است.
ابعاد کوچک سامانه، به‌تنهایی مجوز حذف همهٔ اثرهای اینرسی نیست.
نقش دیواره‌ها نیز باید همراه مقیاس‌های جریان بررسی شود.
\cite{SquiresQuake2005}

در سیال نیوتنی، تنش لزج با نرخ تغییرشکل رابطهٔ خطی دارد.
برای جریان تراکم‌ناپذیر، واگرایی میدان سرعت صفر است.
روی دیوارهٔ ثابت، شرط عدم لغزش سرعت سیال را صفر می‌کند.
\cite{SquiresQuake2005}

برای مقایسهٔ مدل‌ها، طول و سرعت مرجع تعریف می‌شوند.
این مقیاس‌ها به‌ترتیب با $L_0$ و $U_0$ نمایش داده می‌شوند.
مکان، زمان، سرعت و فشار بی‌بعد چنین تعریف می‌شوند:
\begin{equation}
\begin{aligned}
\mathbf r^*&=\frac{\mathbf r}{L_0},&
t^*&=\frac{tU_0}{L_0},\\
\mathbf U^*&=\frac{\mathbf U}{U_0},&
p^*&=\frac{p-p_{\mathrm{ref}}}{\rho_f U_0^2}.
\end{aligned}
\label{eq:lit-reference-scales}
\end{equation}
ستاره، کمیت بی‌بعد را مشخص می‌کند.
$\rho_f$ چگالی سیال و $p_{\mathrm{ref}}$ فشار مرجع است.
مقیاس فشار اینجا، از نوع دینامیکی انتخاب شده است.
این قرارداد، نمونه‌ای توضیحی برای مرور حاضر است.
مقیاس‌های واقعی هر مدل، در فصل سوم مشخص می‌شوند.
اگر مقیاس دو جهت متفاوت باشد، ضرایب مشتق نیز تغییر می‌کنند.
انتخاب فشار مرجع و شرایط خروجی، باید با صورت‌بندی سازگار باشند.

عدد رینولدز، نسبت اثرهای اینرسی و لزجت را بیان می‌کند.
با گرانروی دینامیکی $\mu$، این عدد برابر است با
\begin{equation}
Re=\frac{\rho_f U_0L_0}{\mu}.
\label{eq:lit-reynolds-number}
\end{equation}
بروس، بی‌بعدسازی را مبنای بررسی تقریب جریان استوکس قرار داده است.
حذف جملهٔ همرفتی، به کوچک‌بودن عدد رینولدز وابسته است.
پایابودن و آرام‌بودن جریان نیز باید جداگانه بررسی شوند.
بنابراین، فرض جریان آرام همواره معادل تقریب استوکس نیست.
\cite{Bruus2014Governing}

بروس، جریان پوازوی را نمونه‌ای بنیادی از جریان فشاری معرفی می‌کند.
در حالت توسعه‌یافته، پروفایل محوری در امتداد کانال تغییر نمی‌کند.
برای دو صفحهٔ موازی، پروفایل سرعت سهموی است.
این تقریب، مبنای مدل دوبعدی کانال مستطیلی قرار می‌گیرد.
مقطع مستطیلی سه‌بعدی، به اثر دیواره‌های جانبی نیز وابسته است.
ازاین‌رو، مدل صفحه‌ای باید از مقطع سه‌بعدی متمایز بماند.
\cite{Bruus2014Governing}

پروفایل تحمیلی ورودی، ممکن است با پروفایل توسعه‌یافته متفاوت باشد.
در این وضعیت، ناحیهٔ ورودی بخشی از پاسخ جریان است.
بروس، وابستگی توسعهٔ جریان به رینولدز و ابعاد کانال را بررسی می‌کند.
این موضوع، اهمیت شرایط مرزی را در مقایسهٔ مدل‌ها نشان می‌دهد.
\cite{Bruus2014Governing}
در مرور عملگرهای جریان، هندسه و شرایط مرزی هم‌زمان بررسی می‌شوند.
تعمیم به پروفایل تازه، تنها در دامنهٔ تعریف‌شده معنا دارد.
صورت‌بندی دقیق جریان‌های پژوهش، در فصل سوم ارائه می‌شود.

\زیرقسمت{حرکت لاگرانژی و مدل ذره}
\label{sec:lit-lagrangian-particle}

در توصیف لاگرانژی، موقعیت ذره در طول زمان دنبال می‌شود.
ذرهٔ ردیاب ایده‌آل، با سرعت محلی سیال حرکت می‌کند.
ذرهٔ اینرسی‌دار، می‌تواند سرعتی متفاوت از سیال داشته باشد.
در نتیجه، مسیر آن الزاماً با مسیر ردیاب یکسان نیست.
\cite{ProvencherLanglois2015}

مکسی و رایلی، نیروهای وارد بر کرهٔ کوچک را بازبینی کردند.
صورت‌بندی آنان، اثر تغییرات مکانی جریان را نیز لحاظ می‌کند.
این پژوهش، مرجعی بنیادی برای مدل‌سازی حرکت ذرات است.
\cite{MaxeyRiley1983}
بااین‌حال، استفاده از مدل کاهش‌یافته به توجیه فرض‌ها نیاز دارد.
لانگلوا و همکاران، اهمیت جملهٔ حافظه را در پاسخ گذرا بررسی کردند.
در چارچوب بررسی‌شده، حافظه می‌تواند آهنگ آرامش را تغییر دهد.
بنابراین، حذف آن برای همهٔ شرایط قابل‌تعمیم نیست.
\cite{ProvencherLanglois2015}

برای معرفی سرعت لغزش، سرعت ذره نیز بی‌بعد می‌شود.
با تعریف $\mathbf v_p^*=\mathbf v_p/U_0$، روابط زیر برقرارند:
\begin{equation}
\begin{aligned}
\mathbf w^*&=\mathbf U^*(\mathbf r_p^*,t^*)-\mathbf v_p^*,\\
d_p^*&=\frac{d_p}{L_0},\\
Re_p&=Re\,d_p^*\|\mathbf w^*\|.
\end{aligned}
\label{eq:lit-slip-particle-reynolds}
\end{equation}
در این روابط، $d_p$ قطر ذره است.
عدد رینولدز ذره، بر سرعت لغزش و قطر استوار است.
این عدد، با رینولدز مبتنی بر مقیاس کانال تفاوت دارد.
\cite{MorsiAlexander1972}

مورسی و الکساندر، درگ وابسته به رینولدز لحظه‌ای را بررسی کردند.
آنان، حرکت ذره را در چند میدان جریان مطالعه کردند.
این رویکرد، اهمیت تغییر درگ طی حرکت را نشان می‌دهد.
بنابراین، انتخاب قانون درگ باید با دامنهٔ لغزش سازگار باشد.
\cite{MorsiAlexander1972}
این مطالعه، مرجع مقایسهٔ مدل‌های درگ در مرور حاضر است.
قانون درگ منتخب هر مسئله، در فصل سوم مشخص می‌شود.

عدد استوکس، زمان پاسخ ذره را با زمان مرجع مقایسه می‌کند.
برای کرهٔ دارای درگ استوکس، تعریف زیر استفاده می‌شود:
\begin{equation}
St=\frac{\tau_p U_0}{L_0}
=\frac{\rho_p}{\rho_f}\frac{Re\,(d_p^*)^2}{18}.
\label{eq:lit-stokes-number}
\end{equation}
$\rho_p$ چگالی ذره و $\tau_p$ زمان پاسخ استوکس است.
این تعریف، به فرض کرهٔ منفرد و درگ خطی وابسته است.
این زمان پاسخ، بر جرم خود ذره استوار است.
\cite{ProvencherLanglois2015,Furlani2006Transport}
تغییر قانون درگ، می‌تواند زمان پاسخ مؤثر را تغییر دهد.
عدد استوکس، همراه فرض‌های تعریفش تفسیر می‌شود.

در کوپل یک‌طرفه، میدان حامل مستقل از حرکت ذرات محاسبه می‌شود.
فورلانی، بازخورد ذرات بر جریان را برای سوسپانسیون رقیق حذف کرد.
\cite{Furlani2006Transport}
این فرض، با حذف اینرسی ذره تفاوت دارد.
ذره می‌تواند اینرسی داشته باشد و بازخوردش بر جریان حذف شود.
شرایط اولیه نیز در مدل اینرسی‌دار، شامل مکان و سرعت هستند.
\cite{ProvencherLanglois2015}
پس مقایسهٔ بازسازی‌ها، به ثبت نوع ذره و داده‌ها نیاز دارد.
این اطلاعات، همراه تعداد مسیرها در بخش‌های بعدی بررسی می‌شوند.

\زیرقسمت{حرکت ذرات مغناطیس‌پذیر}
\label{sec:lit-magnetic-particle}

تحریک مغناطیسی، امکان تغییر حرکت ذره در جریان را فراهم می‌کند.
فورلانی، انتقال مغناطیسی را با مدل میدان و درگ بررسی کرد.
\cite{Furlani2006Transport}

فورلانی، مغناطش خطی تا آستانهٔ اشباع را در نظر گرفت.
ناهمگنی فضایی میدان نیز در تعیین نیروی انتقالی نقش دارد.
\cite{Furlani2006Transport}
بنابراین، دامنهٔ فرمان منبع و میدان فضایی باید متمایز بمانند.
تغییر فرمان، از طریق مدل میدان بر نیرو اثر می‌گذارد.

انتخاب مدل حرکت، به زمان پاسخ ذره وابسته است.
برای تحریک متغیر، زمان تغییر فرمان نیز اهمیت دارد.
در تقریب بیش‌میرا، آرامش ذره نسبتاً سریع فرض می‌شود.
در این وضعیت، موازنهٔ نیروها سرعت ذره را تعیین می‌کند.
فورلانی، باوجود کوچک‌بودن اینرسی، آن را در معادله حفظ کرد.
\cite{Furlani2006Transport}

حذف حرکت براونی، نیازمند بررسی مقیاس‌های انتقال است.
اسکوایرز و کوئیک، رقابت جابه‌جایی و پخش را بررسی کرده‌اند.
\cite{SquiresQuake2005}
بنابراین، یک آستانهٔ اندازه برای همهٔ شرایط کافی نیست.
همچنین، نزدیکی دیواره باید هنگام انتخاب مدل درگ بررسی شود.
اهمیت مرزها، از ویژگی‌های اصلی جریان در ابعاد کوچک است.
\cite{SquiresQuake2005}

این مبانی، مقایسهٔ روش‌های یادگیری و بازسازی را هدایت می‌کنند.
برای هر روش، داده‌ها و فرض‌های فیزیکی هم‌زمان ثبت می‌شوند.
شرایط اعتبار مدل‌های اینرسی‌دار و بیش‌میرا نیز جدا بررسی خواهند شد.
مدل میدان و مقیاس‌های مطالعهٔ مغناطیسی، در فصل سوم مشخص می‌شوند.

\قسمت{شبکه‌های عصبی مبتنی بر فیزیک}
\label{sec:lit-pinn}

شبکهٔ عصبی، نمایش پارامتری پاسخ را فراهم می‌کند.
معادلات حاکم نیز می‌توانند در تعیین این پارامترها نقش داشته باشند.
برای توضیح این ارتباط، ابتدا ساختار شبکه و مشتق‌گیری مرور می‌شوند.
سپس، اصول شبکهٔ عصبی مبتنی بر فیزیک معرفی می‌شود.

\زیرقسمت{مبانی شبکهٔ پیشخور و مشتق‌گیری خودکار}
\label{sec:lit-feedforward-autodiff}

گودفلو و همکاران، شبکهٔ پیشخور را ترکیبی از نگاشت‌ها معرفی می‌کنند.
اطلاعات از ورودی، به لایه‌های میانی و خروجی منتقل می‌شود.
در شبکهٔ تمام‌متصل، هر لایه از خروجی لایهٔ قبلی استفاده می‌کند.
وزن‌ها و بایاس‌ها، پارامترهای آموزش‌پذیر این نگاشت هستند.
\cite{Goodfellow2016}

بردار ورودی بی‌بعد با $\mathbf z^*=(\mathbf x^*,t^*)$ نمایش داده می‌شود.
برای مسئلهٔ پایا، مؤلفهٔ زمان از ورودی حذف می‌شود.
یک شبکهٔ پیشخور با $K$ لایهٔ پنهان، چنین نوشته می‌شود:
\begin{equation}
\begin{aligned}
\mathbf h^{(0)}&=\mathbf z^*,\\
\mathbf h^{(\ell)}
&=\sigma_\ell\!\left(
\mathbf W^{(\ell)}\mathbf h^{(\ell-1)}+\mathbf b^{(\ell)}
\right),\qquad \ell=1,\ldots,K,\\
\mathbf N_\theta(\mathbf z^*)
&=\mathbf W^{(o)}\mathbf h^{(K)}+\mathbf b^{(o)}.
\end{aligned}
\label{eq:lit-feedforward-network}
\end{equation}
$\mathbf W^{(\ell)}$ و $\mathbf b^{(\ell)}$ وزن و بایاس لایه هستند.
$\sigma_\ell$ تابع فعال‌ساز و $\theta$ مجموعهٔ پارامترهای شبکه است.
در این نمایش، لایهٔ خروجی خطی انتخاب شده است.
تعداد لایه‌های پنهان و اندازهٔ آن‌ها، عمق و عرض را مشخص می‌کنند.
\cite{Goodfellow2016}

تابع فعال‌ساز، امکان نمایش روابط غیرخطی را فراهم می‌کند.
انتخاب آن، باید با مشتق‌های موردنیاز صورت‌بندی سازگار باشد.
برای مشتق دوم کلاسیک، پاسخ آزمایشی باید دوبار مشتق‌پذیر باشد.
برای مثال، تانژانت هذلولوی مشتق‌های مرتبهٔ بالاتر دارد.
لاگاریس و همکاران، مشتق‌های شبکه را در حل معادلات استفاده کردند.
در روش آنان، مشتق‌پذیری نمایش پاسخ بخشی از ساخت مدل است.
\cite{Lagaris1998}

بایدین و همکاران، روش‌های محاسبهٔ مشتق را دسته‌بندی کرده‌اند.
مشتق‌گیری خودکار، قاعدهٔ زنجیره‌ای را روی عملیات برنامه اعمال می‌کند.
این روش، مقادیر مشتق را هنگام اجرای محاسبات به دست می‌آورد.
در نتیجه، به انتخاب گام تفاضل محدود وابسته نیست.
بااین‌حال، خطاهای محاسبات ممیز شناور همچنان مطرح‌اند.
\cite{Baydin2018}

در کاربرد فیزیکی، مشتق نسبت به مختصات و پارامترها لازم است.
مشتق‌های مکانی و زمانی، باقیماندهٔ معادله را می‌سازند.
گرادیان نسبت به $\theta$، برای به‌روزرسانی پارامترهای شبکه استفاده می‌شود.
پس‌انتشار خطا، با انباشت معکوس مشتق‌ها ارتباط دارد.
این دو کاربرد، از یک گراف محاسباتی مشتق‌پذیر استفاده می‌کنند.
\cite{Baydin2018}

مشتق محاسبه‌شده، مشتق تابع تقریب‌زده‌شدهٔ شبکه است.
دقت مشتق‌گیری، به‌تنهایی دقت پاسخ فیزیکی را تضمین نمی‌کند.
خطای نمایش پاسخ باید با ارزیابی مستقل بررسی شود.
این تمایز، در تفسیر خطای مدل‌های جریان اهمیت دارد.

\زیرقسمت{اصول \lr{PINN} و اعمال قیود فیزیکی}
\label{sec:lit-pinn-principles}

استفاده از شبکه برای حل معادلات، پیشینه‌ای پیش از \lr{PINN} دارد.
لاگاریس و همکاران، پاسخ آزمایشی را از دو بخش ساختند.
بخش معلوم، شرایط اولیه یا مرزی را تأمین می‌کرد.
بخش آموزش‌پذیر، بدون نقض آن شرایط تغییر می‌کرد.
سپس، خطای معادله در نقاط انتخاب‌شده کمینه می‌شد.
\cite{Lagaris1998}

رئیسی و همکاران، چارچوب \lr{PINN} را برای حل و شناسایی توسعه دادند.
در این چارچوب، داده‌ها و معادلات در آموزش مشارکت می‌کنند.
باقیماندهٔ فیزیکی از همان شبکهٔ نمایش پاسخ محاسبه می‌شود.
مشتق‌های لازم نیز با مشتق‌گیری خودکار به دست می‌آیند.
\cite{Raissi2019PINN}

پاسخ تقریب‌زده‌شده با $q_\theta^*(\mathbf z^*)$ نمایش داده می‌شود.
در نمایش اسکالر، نماد $N_\theta$ برای خروجی استفاده می‌شود.
در حالت بدون تبدیل خروجی، $q_\theta^*=N_\theta$ است.
برای بیان عمومی، باقیمانده‌های بی‌بعد زیر تعریف می‌شوند:
\begin{equation}
\begin{aligned}
r_\theta^*(\mathbf z^*)
&=\mathcal N_{\boldsymbol\lambda^*}^{*}
 [q_\theta^*](\mathbf z^*)-f^*(\mathbf z^*),\\
b_\theta^*(\mathbf z_b^*)
&=\mathcal B^{*}[q_\theta^*](\mathbf z_b^*)-g^*(\mathbf z_b^*),\\
i_\theta^*(\mathbf x_0^*)
&=q_\theta^*(\mathbf x_0^*,0)-q_0^*(\mathbf x_0^*).
\end{aligned}
\label{eq:lit-pinn-residuals}
\end{equation}
$\mathcal N^{*}$ عملگر معادله و $\boldsymbol\lambda^*$ ضرایب فیزیکی بی‌بعد هستند.
این عملگر، می‌تواند مشتق‌های مکانی و زمانی را شامل شود.
$\mathcal B^{*}$ نوع شرط مرزی را مشخص می‌کند.
$f^*$، $g^*$ و $q_0^*$ داده‌های تجویزی مسئله هستند.
نمادهای $r_\theta^*$، $b_\theta^*$ و $i_\theta^*$، سه باقیمانده را نشان می‌دهند.
این نمایش، برای دستگاه معادلات نیز قابل‌گسترش است.
در مسئلهٔ پایا، باقیماندهٔ شرط اولیه حذف می‌شود.

صورت جریمه‌ای هزینه، از مجموع خطاهای نمونه‌برداری‌شده ساخته می‌شود.
برای بیان عمومی، چهار مؤلفه در نظر می‌گیریم:
\begin{equation}
\begin{aligned}
\mathcal L
&=\omega_r\mathcal L_r+\omega_b\mathcal L_b
 +\omega_0\mathcal L_0+\omega_d\mathcal L_d,\\
\mathcal L_r
&=\frac{1}{N_r}\sum_{j=1}^{N_r}
 \|r_\theta^*(\mathbf z_{r,j}^*)\|^2,\\
\mathcal L_b
&=\frac{1}{N_b}\sum_{j=1}^{N_b}
 \|b_\theta^*(\mathbf z_{b,j}^*)\|^2,\\
\mathcal L_0
&=\frac{1}{N_0}\sum_{j=1}^{N_0}
 \|i_\theta^*(\mathbf x_{0,j}^*)\|^2,\\
\mathcal L_d
&=\frac{1}{N_d}\sum_{j=1}^{N_d}
 \|\mathcal H^*[q_\theta^*](\mathbf z_{d,j}^*)-d_j^*\|^2.
\end{aligned}
\label{eq:lit-pinn-loss}
\end{equation}
$N_r$، $N_b$، $N_0$ و $N_d$ تعداد نمونه‌های هر گروه هستند.
$\mathcal H^*$ عملگر مشاهده و $d_j^*$ دادهٔ مشاهده‌شدهٔ بی‌بعد است.
این عملگر، می‌تواند فقط مؤلفه‌های اندازه‌گیری‌شده را انتخاب کند.
ضرایب $\omega$، وزن نسبی اجزای هزینه هستند.
مؤلفه‌های نامرتبط، همراه تعداد نمونهٔ متناظر از تعریف حذف می‌شوند.
بنابراین، تعداد نمونهٔ صفر در مخرج وارد نمی‌شود.

رئیسی و همکاران، باقیمانده و داده را در یک هزینه ترکیب کردند.
نمایش بالا، این ایده را با گروه‌بندی صریح‌تر بیان می‌کند.
نقاط فیزیکی، لزوماً مقادیر اندازه‌گیری‌شدهٔ پاسخ را ندارند.
در آن نقاط، سازگاری پاسخ با معادله ارزیابی می‌شود.
\cite{Raissi2019PINN}

شرایط مرزی می‌توانند در ساختار پاسخ نیز اعمال شوند.
سوکومار و سریواستاوا، این مسیر را با توابع فاصله توسعه دادند.
برای نمونهٔ اسکالر دیریکله، نمایش زیر را در نظر می‌گیریم:
\begin{equation}
q_\theta^*(\mathbf z^*)
=A^*(\mathbf z^*)+M^*(\mathbf z^*)N_\theta(\mathbf z^*).
\label{eq:lit-pinn-hard-dirichlet}
\end{equation}
روی مرز دیریکله، $A^*=g^*$ و $M^*=0$ انتخاب می‌شوند.
در نتیجه، مقدار مرزی مستقل از پارامترهای شبکه برقرار می‌ماند.
برای بخش دقیقاً اعمال‌شده، جریمهٔ مرزی ضروری نیست.
توابع سازنده باید مشتق‌پذیری لازم را داشته باشند.
صفرشدن $M^*$، به‌تنهایی شرایط مشتقی را تضمین نمی‌کند.
برای نویمان یا رابین، ساختار متناسب با همان شرط لازم است.
\cite{SukumarSrivastava2022}

در روش جریمه‌ای، قیود از طریق کاهش هزینه تقریب زده می‌شوند.
در روش ساختاری، قید منتخب در نمایش پاسخ برقرار می‌شود.
پس، دقت قید و دقت معادله باید جداگانه ارزیابی شوند.
اعمال دقیق مرز، به‌تنهایی حل دقیق معادله را تضمین نمی‌کند.
\cite{SukumarSrivastava2022}

در مسئلهٔ مستقیم، ضرایب فیزیکی معلوم و پاسخ مجهول است.
در مسئلهٔ معکوس، بعضی ضرایب نیز با داده‌ها تخمین زده می‌شوند.
رئیسی و همکاران، هر دو کاربرد را بررسی کرده‌اند.
\cite{Raissi2019PINN}
بااین‌حال، امکان تخمین به اطلاعات و فرض‌های مسئله وابسته است.
جزئیات شناسایی‌پذیری، در مرور مسائل معکوس بررسی خواهد شد.
در ادامه، آموزش و ارزیابی این شبکه‌ها مرور می‌شوند.

\زیرقسمت{آموزش، نمونه‌برداری و اعتبارسنجی}
\label{sec:lit-pinn-training-validation}

عملکرد \lr{PINN}، به ساخت تابع هزینه و فرایند آموزش وابسته است.
مطالعات مرتبط، چند منشأ متفاوت برای دشواری آموزش گزارش کرده‌اند.
نامتوازنی گرادیان‌ها، توزیع نقاط و بدشرطی، از مهم‌ترین موارد هستند.
بنابراین، مقایسهٔ روش‌ها باید تنظیمات آموزش را نیز دربرگیرد.

وانگ و همکاران، نامتوازنی گرادیان‌های اجزای هزینه را بررسی کردند.
در مسائل بررسی‌شده، این وضعیت یادگیری هم‌زمان قیود را دشوار می‌کرد.
آنان، روشی مبتنی بر آمار گرادیان برای تعدیل آموزش پیشنهاد کردند.
این مطالعه، اهمیت توازن میان داده و قیود فیزیکی را نشان می‌دهد.
بااین‌حال، ضرایب مناسب هزینه به صورت‌بندی مسئله وابسته‌اند.
\cite{WangTengPerdikaris2021}

مقیاس‌بندی، پیش از انتخاب وزن‌های هزینه نیز اهمیت دارد.
برای نمونه، بازهٔ مختصات بی‌بعد می‌تواند به بازه‌ای استاندارد نگاشت شود.
نگاشت نمونه و رابطهٔ مشتق آن، به‌صورت زیر هستند:
\begin{equation}
\widetilde{x}
=2\frac{x^*-a^*}{b^*-a^*}-1,
\qquad
\frac{\partial}{\partial x^*}
=\frac{2}{b^*-a^*}\frac{\partial}{\partial\widetilde{x}}.
\label{eq:lit-pinn-input-scaling}
\end{equation}
در این رابطه، $a^*<b^*$ دو انتهای بازهٔ بی‌بعد هستند.
پس، تغییر ورودی باید با تبدیل مشتق‌ها همراه باشد.
مشتق‌گیری از کل نگاشت، این عامل را در محاسبات وارد می‌کند.
این تبدیل، به‌تنهایی توازن گرادیان‌های هزینه را تضمین نمی‌کند.

انتخاب نقاط فیزیکی، سازوکار دیگری برای بهبود آموزش است.
وو و همکاران، روش‌های ثابت، بازنمونه‌برداری و تطبیقی را مقایسه کردند.
دو روش پیشنهادی آنان، \lr{RAD} و \lr{RAR-D} نام دارند.
در اولی، نقاط براساس توزیع وابسته به باقیمانده جایگزین می‌شوند.
در دومی، نقاط تازه به مجموعهٔ قبلی افزوده می‌شوند.
نتایج مطالعه، مزیت این روش‌ها را در مسائل بررسی‌شده نشان می‌دهد.
\cite{Wu2023Sampling}

این روش‌ها، توزیع تلاش محاسباتی را در دامنه تغییر می‌دهند.
باقیماندهٔ بزرگ، نشانه‌ای برای انتخاب نقاط دشوار فراهم می‌کند.
بااین‌حال، مقدار باقیمانده همیشه برابر خطای پاسخ نیست.
بررسی نقاط مستقل، برای ارزیابی پوشش کل دامنه لازم است.
همچنین، هزینهٔ ارزیابی نقاط نامزد باید در مقایسه لحاظ شود.

راثور و همکاران، دشواری آموزش را از منظر بدشرطی تحلیل کردند.
آنان، نقش عملگرهای دیفرانسیلی را در شکل هزینه بررسی کردند.
در آزمایش‌های آنان، ترکیب \lr{Adam} و \lr{L-BFGS} عملکرد بهتری داشت.
مقایسه، اجرای جداگانهٔ این دو بهینه‌ساز را نیز شامل می‌شد.
آنان، روش مرتبهٔ دوم \lr{NysNewton-CG} را نیز معرفی کردند.
این یافته‌ها، اهمیت انتخاب بهینه‌ساز را در کنار نمونه‌برداری نشان می‌دهند.
\cite{Rathore2024}

این سه مسیر بهبود، نقش‌های متفاوتی دارند.
توازن هزینه، سهم قیود را در به‌روزرسانی پارامترها تنظیم می‌کند.
نمونه‌برداری، مکان ارزیابی معادلات را تغییر می‌دهد.
بهینه‌ساز، مسیر جست‌وجو در فضای پارامترها را تعیین می‌کند.
بنابراین، بهبود یک عامل، کفایت دو عامل دیگر را ثابت نمی‌کند.
برتری گزارش‌شده نیز باید در محدودهٔ آزمایش‌های همان مطالعه تفسیر شود.

راثور و همکاران، هزینهٔ نهایی را همراه خطای پاسخ بررسی کردند.
این مقایسه، ضرورت گزارش دقت فیزیکی کنار هزینه را روشن می‌کند.
\cite{Rathore2024}
برای نمونه، خطای نسبی سرعت روی نقاط مستقل چنین تعریف می‌شود:
\begin{equation}
E_U=
\left[
\frac{\displaystyle\sum_{j=1}^{N_e}
\|\mathbf U_\theta^*(\mathbf x_j^*,t_j^*)
 -\mathbf U_{\mathrm{ref}}^*(\mathbf x_j^*,t_j^*)\|^2}
{\displaystyle\sum_{j=1}^{N_e}
\|\mathbf U_{\mathrm{ref}}^*(\mathbf x_j^*,t_j^*)\|^2}
\right]^{1/2}.
\label{eq:lit-pinn-independent-error}
\end{equation}
در این رابطه، $N_e$ تعداد نقاط ارزیابی مستقل است.
همهٔ سرعت‌ها، با مقیاس مرجع یکسان بی‌بعد شده‌اند.
این تعریف، به غیرصفر بودن مخرج نیاز دارد.
برای مؤلفه‌های بسیار کوچک، خطای مطلق نیز باید گزارش شود.
فشار، سرعت و باقیماندهٔ مومنتوم نیز جداگانه بررسی می‌شوند.
مقایسهٔ فشار، به قرارداد مرجع سازگار نیز نیاز دارد.

در مقایسهٔ آموزش، بودجهٔ محاسباتی باید روشن باشد.
تعداد نقاط، مراحل، سخت‌افزار و زمان آموزش باید ثبت شوند.
نمونه‌های استفاده‌شده برای انتخاب مدل، مجموعهٔ آزمون مستقل محسوب نمی‌شوند.
تکرار اجرا با مقداردهی اولیهٔ متفاوت نیز پایداری را روشن‌تر می‌کند.
این موارد، مبنای بررسی تنظیمات شبکه‌های پژوهش حاضر قرار می‌گیرند.

\زیرقسمت{کاربرد در جریان و انتقال ذرات}
\label{sec:lit-pinn-flow-particle-applications}

کاربردهای جریان، از حل مستقیم تا بازسازی داده‌های آزمایشگاهی گسترده‌اند.
در بررسی این مطالعات، نقش داده و مدل ذره اهمیت دارد.
ذرات می‌توانند ابزار اندازه‌گیری جریان یا موضوع پیش‌بینی حرکت باشند.
این دو نقش، نیازهای داده‌ای متفاوتی ایجاد می‌کنند.

خادمی و همکاران، روش \lr{DG-PINN} را برای جریان ناپایا توسعه دادند.
این روش، دامنهٔ زمانی را به زیردامنه‌های هم‌پوشان تقسیم می‌کند.
هر زیردامنه، شبکهٔ مستقل و پارامترهای خودش را دارد.
انتقال اطلاعات، با سازوکاری وابسته به کیفیت تقریب کنترل می‌شود.
پژوهش، جریان کانالی آشفتهٔ سه‌بعدی را نیز بررسی می‌کند.
\cite{Khademi2025}

در این روش، مقدار پاسخ و گرادیان آن اهمیت دارند.
این انتخاب، پیوستگی پاسخ زیردامنه‌ها را هدف می‌گیرد.
بااین‌حال، انتقال اطلاعات به کیفیت تقریب‌های تولیدشده وابسته می‌ماند.
این رویکرد، با آموزش عملگر برای ورودی‌های تابعی متفاوت یکسان نیست.
بنابراین، نوع تعمیم باید همراه ساختار آموزش مشخص شود.
\cite{Khademi2025}

اشتاینفورت و همکاران، بازسازی میدان از داده‌های رهگیری را بررسی کردند.
در این مطالعه، سرعت‌های حاصل از رهگیری وارد شبکه می‌شوند.
نوع مشاهده و ارزیابی آن، در بخش \ref{sec:lit-timed-flow-inverse} بررسی می‌شود.
\cite{Steinfurth2024}

وان و همکاران، مدل \lr{TrajectoryFlowNet} را برای پیوند مسیر و میدان معرفی کردند.
فرض ردیاب و اطلاعات سرعت، در تفسیر این روش اهمیت دارند.
مقایسهٔ تفصیلی، در بخش \ref{sec:lit-timed-flow-inverse} ارائه می‌شود.
\cite{Wan2026TrajectoryFlowNet}

این مطالعات، سه جهت متفاوت توسعه را نشان می‌دهند.
تقسیم دامنه، دشواری حل جریان ناپایا را هدف می‌گیرد.
همگون‌سازی داده، مشاهدات پراکنده را به میدان متراکم تبدیل می‌کند.
پیوند لاگرانژی–اویلری، سازگاری میدان و حرکت را دنبال می‌کند.
ارزیابی هر جهت، باید با هدف و فرض‌های همان مسئله انجام شود.

برای پژوهش حاضر، چند تمایز در مقایسه‌ها حفظ می‌شود.
خطای میدان و خطای مسیر، دو معیار جدا هستند.
مدل درگ، اطلاعات اولیه و نوع مشاهده نیز ثبت می‌شوند.
فرض ردیاب، بدون بررسی به ذرهٔ اینرسی‌دار تعمیم داده نمی‌شود.
جزئیات مقایسهٔ بازسازی‌ها، در بخش مسائل معکوس ارائه خواهد شد.
در بخش بعد، یادگیری پاسخ خانواده‌ای از مسائل بررسی می‌شود.

\قسمت{یادگیری عملگر و شبکه‌های \lr{DeepONet}}
\label{sec:lit-operator-learning}

\زیرقسمت{مفهوم یادگیری عملگر}
\label{sec:lit-operator-concept}

در بسیاری از مسائل، پاسخ به یک ورودی تابعی وابسته است.
این ورودی می‌تواند شرط اولیه، پروفایل مرزی یا تابع تحریک باشد.
تغییر ورودی، معمولاً پاسخ متناظر دیگری ایجاد می‌کند.
یادگیری عملگر، نگاشت میان این ورودی‌ها و پاسخ‌ها را تقریب می‌زند.
بنابراین، هدف از تقریب یک پاسخ منفرد فراتر می‌رود.
\cite{Kovachki2023NeuralOperator}

برای بیان روابط، ابتدا مقیاس‌های مرجع مثبت انتخاب می‌شوند.
طول، سرعت، دامنهٔ ورودی و دامنهٔ پاسخ، مقیاس‌های متناسب خود را دارند.
در مسائل انتقالی، زمان مرجع از طول و سرعت ساخته می‌شود.
تعریف متغیرهای بی‌بعد به‌صورت زیر است:
\begin{equation}
\mathbf{x}^{*}=\frac{\mathbf{x}}{L_0},
\qquad T_0=\frac{L_0}{U_0},
\qquad t^{*}=\frac{t}{T_0},
\qquad a^{*}=\frac{a}{A_0},
\qquad w^{*}=\frac{w}{W_0}.
\label{eq:lit-operator-scales}
\end{equation}
در خروجی برداری، هر کمیت مقیاس متناسب خود را دارد.
این بی‌بعدسازی، با نرمال‌سازی آماری ورودی شبکه تفاوت دارد.
هیچ‌کدام، به‌تنهایی شرایط مرزی یا هندسهٔ مسئله را تغییر نمی‌دهند.

ورودی تابعی بی‌بعد را با $a^{*}$ نمایش می‌دهیم.
بردار $\boldsymbol{\lambda}^{*}$، پارامترهای بی‌بعدِ معلوم را دربرمی‌گیرد.
پاسخ متناظر، تابع $w^{*}$ از مختصات پرس‌وجو است.
برای تعریف نگاشت، پاسخ هر ورودی باید مشخص باشد.
در مسائل چندپاسخی، باید شاخهٔ پاسخ نیز تعیین شود.
در این مرور، خانوادهٔ مسائل دارای پاسخ مشخص فرض می‌شود.
این مختصات، بسته به مسئله مکانی، زمانی یا ترکیبی هستند.
نگاشت موردنظر را می‌توان چنین نوشت:
\begin{equation}
\mathcal{G}^{*}:\mathcal{A}^{*}\times\Lambda^{*}\longrightarrow\mathcal{W}^{*},
\qquad
w^{*}(\mathbf{z}^{*})
=\mathcal{G}^{*}[a^{*},\boldsymbol{\lambda}^{*}](\mathbf{z}^{*}).
\label{eq:lit-operator-map}
\end{equation}
مجموعه‌های $\mathcal{A}^{*}$ و $\mathcal{W}^{*}$، فضاهای مجاز ورودی و پاسخ هستند.
مجموعهٔ $\Lambda^{*}$ نیز دامنهٔ پارامترهای معلوم را مشخص می‌کند.
تابع ورودی با نقطهٔ پرس‌وجو، نقش یکسانی ندارد.
اولی مسئله را تعیین می‌کند؛ دومی محل ارزیابی پاسخ است.
\cite{lu2021deeponet}

در حالت معمول، \lr{PINN} پاسخ یک مسئلهٔ مشخص را تقریب می‌زند.
البته، افزودن پارامترها به \lr{PINN} نیز امکان‌پذیر است.
تمایز اصلی، نوع نگاشت هدف و نحوهٔ نمایش ورودی تابعی است.
در یادگیری عملگر، پارامترهای شبکه میان ورودی‌های مختلف مشترک هستند.
پس از آموزش، ورودی تازه می‌تواند بدون بازآموزی ارزیابی شود.
اعتبار این ارزیابی، به دامنهٔ آموزش و آزمون وابسته است.
\cite{wang2021pideeponet,Kovachki2023NeuralOperator}

در این زمینه، دو معنای تعمیم باید جدا شوند.
ارزیابی در مختصات تازه، تعمیم نسبت به نقطهٔ پرس‌وجو است.
پیش‌بینی پاسخ ورودی تازه، تعمیم نسبت به تابع ورودی است.
موفقیت در اولی، به‌تنهایی دومی را اثبات نمی‌کند.
آزمون درون خانواده و بیرون خانواده نیز باید تفکیک شوند.
برای مثال، تغییر دامنه یا فرکانس تحریک، آزمون‌های متفاوتی می‌سازد.
خانوادهٔ آموزش، دامنهٔ پارامترها و تعریف آزمون باید روشن باشند.

هزینهٔ آموزش، از هزینهٔ ارزیابی پس از آموزش جدا است.
ساخت داده و محاسبهٔ مشتق‌ها، ممکن است آموزش را پرهزینه کنند.
این هزینه، در کاربردهای دارای ارزیابی مکرر اهمیت بیشتری دارد.
صرفه‌جویی نهایی باید همراه دقت و تعداد ارزیابی‌ها سنجیده شود.
خانواده‌های عملگر عصبی، از نمایش‌های طیفی یا انتگرالی نیز استفاده می‌کنند.
مرور حاضر، بر معماری \lr{DeepONet} و پیوند آن با فیزیک تمرکز دارد.
\cite{Kovachki2023NeuralOperator}

\زیرقسمت{معماری شاخه--تنه و نقش حسگرها}
\label{sec:lit-deeponet-architecture}

لو و همکاران، معماری \lr{DeepONet} را برای تقریب عملگر معرفی کردند.
این معماری، دو شبکهٔ شاخه و تنه دارد.
شاخه، نمایش ورودی تابعی را می‌سازد.
تنه، مختصات نقطهٔ پرس‌وجو را دریافت می‌کند.
خروجی نهایی، از ترکیب نمایش این دو شبکه به‌دست می‌آید.
\cite{lu2021deeponet}

در نمایش متداول، تابع ورودی روی چند حسگر نمونه‌برداری می‌شود.
واژهٔ حسگر، اینجا به مکان نمونه‌برداری اشاره دارد.
وجود ابزار اندازه‌گیری فیزیکی، برای این تعریف ضروری نیست.
بردار نمونه‌ها و نمایش دو شبکه چنین تعریف می‌شوند:
\begin{equation}
\begin{aligned}
\mathbf{a}^{*}_{m}
&=\big[a^{*}(s^{*}_{1}),\ldots,a^{*}(s^{*}_{m})\big]^{\mathsf T},\\
\mathbf{b}_{\theta_B}
&=B_{\theta_B}(\mathbf{a}^{*}_{m},\boldsymbol{\lambda}^{*})\in\mathbb{R}^{q},\\
\mathbf{t}_{\theta_T}
&=T_{\theta_T}(\mathbf{z}^{*})\in\mathbb{R}^{q}.
\end{aligned}
\label{eq:lit-branch-trunk-features}
\end{equation}
نماد $m$، تعداد حسگرهای هر ورودی را مشخص می‌کند.
نماد $q$، تعداد ویژگی‌های متناظر دو شبکه است.
پارامترهای اسکالر می‌توانند کنار نمونه‌ها وارد شاخه شوند.
این توسعه، باید در تعریف ورودی مدل صریح باشد.
برای یک خروجی اسکالر، ترکیب متداول چنین است:
\begin{equation}
\widehat{w}^{*}_{\theta}(\mathbf{z}^{*};a^{*},\boldsymbol{\lambda}^{*})
=\sum_{k=1}^{q}
b_{\theta_B,k}(\mathbf{a}^{*}_{m},\boldsymbol{\lambda}^{*})
t_{\theta_T,k}(\mathbf{z}^{*})+c_{\theta}.
\label{eq:lit-deeponet-contraction}
\end{equation}
ضرایب شاخه، با تغییر تابع ورودی تغییر می‌کنند.
توابع تنه، نمایش آموخته‌شده‌ای از وابستگی پاسخ به مختصات هستند.
این توابع، لزوماً متعامد یا دارای تفسیر فیزیکی مستقل نیستند.
تعداد ویژگی‌ها نیز با تعداد نقاط خروجی برابر نیست.
برای خروجی برداری، می‌توان ترکیب‌های متناظر هر مؤلفه را تعریف کرد.
\cite{lu2021deeponet}

در \lr{DeepONet} متداول، تعداد و ترتیب حسگرهای ورودی ثابت هستند.
آزادی انتخاب نقاط خروجی، این محدودیت ورودی را حذف نمی‌کند.
کوواچکی و همکاران، این تفاوت را در تحلیل عملگرهای عصبی بررسی کردند.
بنابراین، خروجی پیوسته با استقلال کامل از گسسته‌سازی یکسان نیست.
\cite{Kovachki2023NeuralOperator}

به‌شرط دسترسی انحصاری به حسگرها، یک محدودیت اطلاعاتی وجود دارد.
دو تابع متفاوت ممکن است روی همهٔ حسگرها یکسان باشند.
در این حالت، شبکه نمایش ورودی یکسانی دریافت می‌کند.
درنتیجه، نمی‌تواند پاسخ‌های متفاوت آن‌ها را از این اطلاعات تشخیص دهد.
افزایش تعداد توابع آموزشی، این ابهام اطلاعاتی را لزوماً رفع نمی‌کند.
افزایش حسگرها نیز جای تنوع توابع آموزشی را نمی‌گیرد.
البته، اعمال قیود می‌تواند مسیر دیگری برای دریافت اطلاعات ورودی بسازد.
مثلاً، یک انتگرال تحلیلی می‌تواند مستقیماً در تبدیل خروجی وارد شود.
در چنین مدلی، همهٔ اطلاعات ورودی از شاخه عبور نمی‌کنند.

بهمنی و همکاران، در سال ۲۰۲۵ این محدودیت گسسته‌سازی را بررسی کردند.
روش آنان، نمونه‌های نامنظم را به نمایش ضرایب پایه‌ها تبدیل می‌کند.
این نمایش، ورودی \lr{DeepONet} را از شبکهٔ نمونه‌برداری اولیه جدا می‌سازد.
آنان، توسعه‌ای مشابه برای نمایش خروجی نیز ارائه کردند.
بااین‌حال، نمونه‌برداری همچنان باید اطلاعات کافی از تابع داشته باشد.
انعطاف حسگرها، به معنای بازیابی اطلاعات حذف‌شده نیست.
\cite{Bahmani2025RINO}

قضیهٔ تقریب همگانی، پشتوانه‌ای برای امکان تقریب عملگر فراهم می‌کند.
صورت کلاسیک، پیوستگی عملگر و فشردگی مجموعهٔ ورودی را فرض می‌کند.
تعداد حسگرها و ظرفیت شبکه، می‌توانند همراه دقت مطلوب تغییر کنند.
این قضیه، موفقیت آموزش یک معماری ثابت را تضمین نمی‌کند.
همچنین، تضمینی عمومی برای پاسخ خارج از خانوادهٔ مفروض نمی‌دهد.
\cite{lu2021deeponet}

برای گزارش آموزش، پنج دستهٔ اطلاعات باید جدا ثبت شوند.
جدول~\ref{tab:lit-operator-counts}، معنای این دسته‌ها را نشان می‌دهد.
تکرار تابع در چند دستهٔ آموزشی، تابع متمایز تازه‌ای نمی‌سازد.
در نمونه‌سازی برخط، ورودی‌ها هنگام آموزش تولید می‌شوند.
در مخزن ثابت، آموزش از مجموعه‌ای ازپیش‌ساخته نمونه می‌گیرد.
هیچ‌کدام، به‌تنهایی معیار درستی یا نادرستی \lr{DeepONet} نیستند.
پوشش خانوادهٔ ورودی و آزمون مستقل، ملاک ارزیابی هستند.

\begin{table}[htbp]
\centering
\caption{کمیت‌های متمایز در گزارش آموزش عملگر}
\label{tab:lit-operator-counts}
\begin{tabular}{p{0.30\textwidth}p{0.59\textwidth}}
\hline
کمیت & معنای مورد گزارش \\
\hline
توابع آموزشی & تعداد ورودی‌های تابعی متمایز و روش تولید آن‌ها \\
حسگرهای ورودی & تعداد و مکان نمونه‌های هر تابع \\
نقاط پرس‌وجو & تعداد نقاط پاسخ یا باقیمانده برای هر ورودی \\
بودجهٔ آموزش & گام‌های بهینه‌ساز و اندازهٔ دستهٔ توابع و نقاط \\
اعتبارسنجی و آزمون & تعداد توابع مستقل و روش جداسازی آن‌ها \\
\hline
\end{tabular}
\end{table}

\زیرقسمت{شبکهٔ \lr{DeepONet} مبتنی بر فیزیک}
\label{sec:lit-physics-informed-deeponet}

وانگ و همکاران، قیود فیزیکی را وارد آموزش \lr{DeepONet} کردند.
در روش آنان، مشتق‌های خروجی با مشتق‌گیری خودکار محاسبه می‌شوند.
سپس، معادلات حاکم و شرایط مسئله در هزینه وارد می‌شوند.
این ساختار، نیاز به پاسخ‌های برچسب‌دار را می‌تواند کاهش دهد.
در مثال‌های مقاله، آموزش بدون پاسخ‌های برچسب‌دار نیز بررسی شده است.
\cite{wang2021pideeponet}

برای بیان عمومی، معادلهٔ بی‌بعد مسئله را با $\mathcal{N}^{*}$ نشان می‌دهیم.
این نماد، عملگر دیفرانسیلیِ معلومِ مسئله است.
عملگر یادگرفته‌شدهٔ $\mathcal{G}^{*}_{\theta}$، نگاشت ورودی به پاسخ را تقریب می‌زند.
این دو عملگر، نقش متفاوتی دارند.
برای ورودی آموزشی $i$، باقیمانده چنین تعریف می‌شود:
\begin{equation}
r^{*}_{\theta,i}(\mathbf{z}^{*})
=\mathcal{N}^{*}\!\left[
\widehat{w}^{*}_{\theta}(\cdot;a_i^{*},\boldsymbol{\lambda}_i^{*});
a_i^{*},\boldsymbol{\lambda}_i^{*}
\right](\mathbf{z}^{*}).
\label{eq:lit-operator-physical-residual}
\end{equation}
با فرض تعداد نقاط یکسان، نمونه‌ای از هزینه چنین است:
\begin{equation}
\begin{aligned}
\mathcal{L}_{\mathrm{phys}}(\theta)
&=\frac{1}{N_fN_r}\sum_{i=1}^{N_f}\sum_{j=1}^{N_r}
\left\|r^{*}_{\theta,i}(\mathbf{z}^{*}_{ij})\right\|_2^2,\\
\mathcal{L}(\theta)
&=\omega_r\mathcal{L}_{\mathrm{phys}}
+\omega_b\mathcal{L}_{\mathrm{BC}}
+\omega_0\mathcal{L}_{\mathrm{IC}}
+\omega_d\mathcal{L}_{\mathrm{data}}.
\end{aligned}
\label{eq:lit-operator-physical-loss}
\end{equation}
نماد $N_f$، تعداد توابع و $N_r$، تعداد نقاط داخلی هرکدام است.
اجزای مرزی و اولیه، ناسازگاری با قیود متناظر را اندازه می‌گیرند.
جزء داده، اختلاف پاسخ شبکه با پاسخ مشاهده‌شده را می‌سنجد.
وزن‌های $\omega$، سهم اجزای هزینه را تعیین می‌کنند.
در نبود برچسب پاسخ، وزن $\omega_d$ صفر قرار می‌گیرد.
در مسئلهٔ پایا، جزء شرط اولیه نیز حذف می‌شود.
روابط بالا، صورت توضیحی هستند؛ تنظیمات واقعی در فصل سوم می‌آیند.

آموزش داده‌محور، از زوج‌های ورودی و پاسخ استفاده می‌کند.
آموزش فیزیکی، پاسخ را با معادلات و قیود محدود می‌سازد.
آموزش ترکیبی، از هر دو نوع اطلاعات بهره می‌گیرد.
نبود برچسب پاسخ، به معنای نبود اطلاعات مسئله نیست.
خانوادهٔ ورودی، ضرایب معادله و شرایط مسئله همچنان لازم هستند.
کاهش هزینهٔ تولید برچسب، الزاماً هزینهٔ آموزش را کاهش نمی‌دهد.
مشتق‌های مرتبهٔ بالاتر می‌توانند محاسبات آموزش را سنگین کنند.
\cite{wang2021pideeponet}

قیود مرزی می‌توانند به‌صورت جریمه یا ساختاری اعمال شوند.
برای شرط دیریکله، تبدیل عمومی زیر قابل استفاده است:
\begin{equation}
\widehat{w}^{*}_{\theta}
=w^{*}_{\mathrm{lift}}+M^{*}\widetilde{w}^{*}_{\theta},
\qquad
w^{*}_{\mathrm{lift}}\big|_{\Gamma_D^{*}}=g_D^{*},
\qquad M^{*}\big|_{\Gamma_D^{*}}=0.
\label{eq:lit-operator-hard-boundary}
\end{equation}
تابع نخست، مقدار مرزی مطلوب را تأمین می‌کند.
تابع $M^{*}$، سهم تصحیح شبکه را روی مرز صفر می‌کند.
این تبدیل، نمونه‌ای عمومی از اعمال دقیق شرایط مرزی است.
ساخت آن، باید با هندسه و همواری موردنیاز سازگار باشد.
\cite{SukumarSrivastava2022}

در جریان دوبعدی تراکم‌ناپذیر، نمایش تابع جریان نیز مفید است.
با مقیاس $U_0L_0$ برای تابع جریان، داریم:
\begin{equation}
u^{*}=\frac{\partial\psi^{*}}{\partial y^{*}},
\qquad
v^{*}=-\frac{\partial\psi^{*}}{\partial x^{*}},
\qquad
\frac{\partial u^{*}}{\partial x^{*}}
+\frac{\partial v^{*}}{\partial y^{*}}=0.
\label{eq:lit-operator-streamfunction}
\end{equation}
پیوستگی، با فرض وجود مشتق‌های مختلط پیوسته، خودکار برقرار می‌شود.
بااین‌حال، مومنتوم و همهٔ شرایط مرزی، خودکار اعمال نمی‌شوند.
انتخاب چنین ساختاری، از افزودن صرف جریمهٔ فیزیکی متمایز است.
معادلات اختصاصیِ مدل‌های این پژوهش، در فصل سوم ارائه می‌شوند.

پژوهش‌های تازه، دشواری آموزش عملگر فیزیکی را نیز بررسی کرده‌اند.
هوانگ و همکاران، در سال ۲۰۲۵ روش \lr{DONIS} را معرفی کردند.
این روش، انتخاب توابع و نقاط باقیمانده را جدا تنظیم می‌کند.
اهمیت نمونه‌ها، با اطلاعاتی از خطای فیزیکی تخمین زده می‌شود.
در مثال‌های مقاله، این انتخاب آموزش را بهبود داده است.
این مطالعه، اهمیت نمونه‌برداری در هر دو فضای ورودی را نشان می‌دهد.
بااین‌حال، بهبود گزارش‌شده، تضمینی برای هر مسئله یا بودجه نیست.
\cite{Huang2025DONIS}

یانگ، در مقالهٔ ۲۰۲۶، تقریب و تعمیم \lr{DeepONet} فیزیکی را تحلیل کرد.
این تحلیل، خطای مشتق‌ها را نیز در نرم‌های سوبولف بررسی می‌کند.
در چارچوب بررسی‌شده، افزایش ظرفیت شاخه نقش برجسته‌ای دارد.
اما افزایش پیچیدگی تنه، همواره سود مشابهی ایجاد نمی‌کند.
این نتیجه، توصیه‌ای مطلق برای همهٔ معماری‌ها و مسائل نیست.
تحلیل همچنین، نمونه‌برداری توابع را از نمونه‌برداری مختصات تفکیک می‌کند.
\cite{Yang2026PIDeepONet}

در نتیجه، ارزیابی باید روی توابع کنارگذاشته‌شده انجام شود.
کنارگذاشتن چند نقطه از همان تابع، برای آزمون عملگر کافی نیست.
خطای پاسخ و باقیماندهٔ فیزیکی نیز باید جدا گزارش شوند.
برای میدان برداری، خطای مؤلفه‌ها می‌تواند رفتارهای متفاوتی نشان دهد.
در پیش‌بینی زنجیره‌ای، انباشت خطا نیز به بررسی مستقل نیاز دارد.
این ملاحظات، مبنای مرور کاربردهای جریان و حرکت در ادامه هستند.

\زیرقسمت{یادگیری جریان برای ورودی‌های متفاوت}
\label{sec:lit-variable-inlet-flow}

در جریان کانالی، تغییر ورودی می‌تواند میدان داخلی را تغییر دهد.
این تغییر، ممکن است مربوط به شدت یا شکل پروفایل باشد.
تغییر ویسکوزیته نیز عدد رینولدز را تغییر می‌دهد.
این حالت‌ها، اطلاعات و دامنه‌های تعمیم متفاوتی ایجاد می‌کنند.
یادگیری عملگر، امکان تقریب پاسخ خانواده‌ای از این مسائل را فراهم می‌کند.
\cite{lu2021deeponet,wang2021pideeponet}

برای نمونه، پروفایل سرعت محوری ورودی را با $g(y)$ نشان می‌دهیم.
ارتفاع کانال مستطیلی، $H$ و ویسکوزیتهٔ سینماتیکی، $\nu$ هستند.
دو قرارداد مقیاس‌بندی، باید از یکدیگر متمایز شوند.
اگر سرعت مرجع ثابت باشد، دامنهٔ ورودی در پروفایل حفظ می‌شود.
اگر میانگین هر نمونه مبنا باشد، پروفایل میانگین واحد دارد.
در قرارداد دوم، تغییر میانگین در عدد رینولدز ظاهر می‌شود.
برای میانگین مثبت، تعریف‌های متناظر چنین هستند:
\begin{equation}
\begin{aligned}
\overline{U}&=\frac{1}{H}\int_0^H g(y)\,\mathrm{d}y>0,
&\eta&=\frac{y}{H},\\
f^{*}(\eta)&=\frac{g(H\eta)}{\overline{U}},
&\int_0^1 f^{*}(\eta)\,\mathrm{d}\eta&=1,\\
\mathbf{x}^{*}&=\frac{\mathbf{x}}{H},
&\mathbf{u}^{*}&=\frac{\mathbf{u}}{\overline{U}},\\
p^{*}&=\frac{p-p_{\mathrm{ref}}}{\rho\overline{U}^{2}},
&Re&=\frac{\overline{U}H}{\nu}.
\end{aligned}
\label{eq:lit-inlet-shape-scaling}
\end{equation}
در این روابط، $\rho$ چگالی و $p_{\mathrm{ref}}$ مرجع فشار است.
انتخاب قرارداد، باید با معادلات بی‌بعد و داده‌های ورودی سازگار باشد.
در نتیجه، حذف دامنهٔ سرعت بدون نگه‌داشتن اطلاعات معادل، مجاز نیست.

با هندسهٔ ثابت، نگاشت توضیحی جریان را می‌توان چنین نوشت:
\begin{equation}
\mathcal{G}^{*}_{\Omega^{*}}
\left[f^{*},Re;\mathcal{B}^{*}_{\mathrm{out}}\right](x^{*},y^{*})
=\left(u^{*},v^{*},p^{*}\right)(x^{*},y^{*}).
\label{eq:lit-inlet-flow-operator}
\end{equation}
نماد $\mathcal{B}^{*}_{\mathrm{out}}$، قرارداد شرط خروجی را مشخص می‌کند.
این قرارداد می‌تواند در تمام خانوادهٔ مسائل ثابت باشد.
در این صورت، ورودی متغیر شبکه محسوب نمی‌شود.
شکل ورودی و رینولدز، به‌تنهایی شرایط کامل مسئله را تعیین نمی‌کنند.
هندسه، قیود دیواره و شرط خروجی نیز باید معلوم باشند.
خانوادهٔ ورودی‌ها باید با این قیود سازگار انتخاب شود.

مطالعات پیشین، دامنه‌های متفاوتی از این نگاشت را بررسی کرده‌اند.
لو و همکاران، در سال ۲۰۲۲ چند عملگر عصبی را مقایسه کردند.
یکی از مثال‌ها، جریان حفرهٔ دوبعدی با درپوش متحرک بود.
داده‌های مرجع این مثال، با روش بولتزمن شبکه‌ای تولید شدند.
مطالعه، تأثیر نمایش ورودی و توسعه‌های معماری را نیز بررسی کرد.
بااین‌حال، این مثال ورودی و خروجی باز ندارد.
بنابراین، تعمیم شرط درپوش با تعمیم ورودی کانال یکسان نیست.
خانوادهٔ قیود بررسی‌شده نیز با پروفایل دلخواه برابر نیست.
\cite{Lu2022FairComparison}

اونیشی و همکاران، در سال ۲۰۲۶ مسئله‌ای نزدیک‌تر بررسی کردند.
آنان، جریان پایای دوبعدی را با \lr{PI-DeepONet} تقریب زدند.
تغییر رینولدز، در جریان پلهٔ پس‌رونده بررسی شد.
تغییر شکل ورودی، در یک کانال جداگانه بررسی شد.
آموزش، بدون برچسب میدان و با قیود فیزیکی انجام گرفت.
برای کانال، ۸۰۰ پروفایل و ۵۰ حسگر ورودی استفاده شد.
پروفایل‌ها، از ترکیب توابع پایهٔ شعاعی ساخته شدند.
این دو تعداد، به‌ترتیب تنوع ورودی و تفکیک نمونه‌برداری را بیان می‌کنند.
\cite{Onishi2026FlowConditions}

مرجع عددی این مطالعه، حل حجم محدود با \lr{OpenFOAM} بود.
شرط خروجی، فشار معلوم و گرادیان نرمالِ صفر سرعت بود.
گرادیان نرمال، مشتق در راستای عمود بر مرز خروجی است.
این شرط، با هر صورت شرط طبیعی خروجی یکسان نیست.
آزمون‌ها، پروفایل‌های سهموی، چندجمله‌ای و جت موضعی را شامل شدند.
افت دقت برای برخی تغییرات دامنه و آزمون‌های برون‌یابی گزارش شد.
پس، حضور فیزیک در آموزش، دامنهٔ تعمیم را نامحدود نمی‌کند.
\cite{Onishi2026FlowConditions}

کروز گونزالس و همکاران، جریان آنوریسم سه‌بعدی را بررسی کردند.
ورودی‌های مدل، شرایط سرعت ورودی و فشار خروجی را دربرمی‌گرفتند.
خروجی‌ها، مؤلفه‌های سرعت و فشار میدان بودند.
آموزش مدل فیزیکی، از دادهٔ حل عددی نیز استفاده می‌کرد.
در خانوادهٔ اصلی، شدت پروفایل سهموی ورودی تغییر می‌کرد.
تغییر شدت، با یادگیری شکل‌های مکانی مستقل تفاوت دارد.
همچنین، هندسهٔ سه‌بعدی این مطالعه با کانال دوبعدی یکسان نیست.
\cite{CruzGonzalez2026Vascular}

روش‌های مبتنی بر فیزیک، به معماری \lr{DeepONet} محدود نیستند.
لی و همکاران، روش \lr{PINO} را بر پایهٔ \lr{FNO} ارائه کردند.
آنان، دادهٔ کم‌تفکیک را با قیود فیزیکی پُرتفکیک ترکیب کردند.
آزمایش‌ها، جریان ناویر--استوکس و انتقال میان رینولدزها را نیز دربرمی‌گرفتند.
در بخشی از روش، پالایش مخصوص هر مسئله انجام می‌شد.
این پالایش، هزینه‌ای افزون بر ارزیابی شبکهٔ ثابت دارد.
پس، مقایسهٔ سرعت باید نقش بازآموزی یا پالایش را روشن کند.
\cite{Li2024PINO}

جدول~\ref{tab:lit-variable-inlet-studies}، تفاوت دامنهٔ این مطالعات را خلاصه می‌کند.
وجود واژهٔ عملگر، برای یکسان‌دانستن آزمون‌های آن‌ها کافی نیست.
باید نوع ورودی، نظارت و مرجع هر مطالعه مشخص باشد.

\begin{table}[htbp]
\centering
\caption{تفاوت دامنهٔ مطالعات یادگیری عملگر جریان}
\label{tab:lit-variable-inlet-studies}
\begin{tabular}{p{0.16\textwidth}p{0.25\textwidth}p{0.23\textwidth}p{0.24\textwidth}}
\hline
مطالعه & تغییر بررسی‌شده & اطلاعات آموزش & تفاوت با مسئلهٔ کانال \\
\hline
لو، ۲۰۲۲ & درپوش و حالات جریان حفره & دادهٔ بولتزمن شبکه‌ای & دامنهٔ بسته و خانوادهٔ محدود قیود \\
اونیشی، ۲۰۲۶ & رینولدز و پروفایل در دو مثال & معادلات و شرایط مرزی & نزدیک‌ترین نمونهٔ مستقیمِ مرورشده \\
کروز گونزالس، ۲۰۲۶ & شدت ورودی و شرایط خروجی & دادهٔ عددی و قیود فیزیکی & هندسهٔ عروقی سه‌بعدی \\
لی، ۲۰۲۴ & خانواده‌های معادلات و رینولدز & داده، فیزیک و پالایش موردی & معماری طیفی و نقش پالایش \\
\hline
\end{tabular}
\end{table}

برای مقایسهٔ منصفانه، تعریف خانوادهٔ پروفایل‌ها اهمیت دارد.
پروفایل سهموی، تخت و نامتقارن، اطلاعات مکانی متفاوتی دارند.
پروفایل نوسانی نیز می‌تواند تفکیک بیشتری از ورودی بخواهد.
در دیوارهٔ بدون لغزش، سازگاری سرعت گوشه‌های ورودی باید بررسی شود.
مثلاً، ورودی تخت ممکن است نزدیک دیواره به هموارسازی نیاز داشته باشد.
این هموارسازی، باید هنگام معرفی خانوادهٔ آموزش مشخص شود.
همچنین، دامنهٔ سرعت و رینولدز باید همراه شکل پروفایل گزارش شوند.

اعتبارسنجی، باید مسئلهٔ عددی یکسانی را مقایسه کند.
طول کانال و شرط خروجی، می‌توانند بر میدان ورودی‌توسعه‌یابنده اثر بگذارند.
بنابراین، مقایسهٔ دو هندسه یا خروجی متفاوت، آزمون دقت مستقیم نیست.
حل مرجع می‌تواند المان محدود، حجم محدود یا روش معتبر دیگری باشد.
سازگاری شرایط و بررسی تفکیک حل مرجع، باید روشن باشند.
فشار نیز باید با قرارداد مرجع یکسان مقایسه شود.
وقتی خروجی فشار را تثبیت می‌کند، جابه‌جایی دلخواه آن مجاز نیست.
در مسائل دارای آزادی ثابت، مرجع مشترک انتخاب می‌شود.

خطای بردار سرعت، همهٔ تفاوت‌های میدان را آشکار نمی‌کند.
اگر سرعت محوری غالب باشد، خطای عرضی ممکن است کم‌رنگ شود.
بنابراین، خطای مؤلفهٔ عرضی باید جدا گزارش شود.
نزدیک صفر بودن این مؤلفه، تفسیر خطای نسبی را دشوار می‌کند.
در این حالت، خطای مطلق بی‌بعد نیز مفید است.
خطای فشار، افت فشار و باقیماندهٔ مومنتوم نیز اطلاعات مکمل دارند.
هیچ‌کدام، به‌تنهایی جای سایر معیارها را نمی‌گیرند.

در نهایت، آزمون باید توابع ورودیِ کنارگذاشته‌شده را شامل شود.
نقاط مکانی تازه از یک پروفایل آموزشی، آزمون محدودتری هستند.
تغییر هم‌زمان پروفایل و رینولدز نیز باید جدا ارزیابی شود.
برای اتصال به حرکت ذره، دقت میدان همچنان اهمیت دارد.
بازتولید یک مسیر، به‌تنهایی درستی کل میدان را اثبات نمی‌کند.
ازاین‌رو، مقایسهٔ میدان و بازپخش مستقل مسیر، مکمل یکدیگر هستند.
این تفکیک، مبنای بررسی مسائل معکوس در ادامهٔ فصل خواهد بود.

\زیرقسمت{یادگیری پاسخ حرکتی برای تحریک زمانی}
\label{sec:lit-temporal-forcing-operators}

در دستگاه‌های تحریک‌شونده، پاسخ به تاریخچهٔ ورودی وابسته است.
این ورودی می‌تواند نیرو، شتاب پایه یا فرمان منبع مغناطیسی باشد.
یادگیری عملگر، نگاشت میان تابع تحریک و پاسخ را تقریب می‌زند.
نمونهٔ آونگ تحریک‌شونده، در پژوهش اولیهٔ \lr{DeepONet} بررسی شده است.
بااین‌حال، شباهت نوع نگاشت، به معنای یکسانی مدل فیزیکی نیست.
\cite{lu2021deeponet}

برای توضیح، دستگاهی با حالت کامل $\mathbf{q}$ درنظر می‌گیریم.
مکان با $L_0$ و سرعت با $U_0$ بی‌بعد می‌شوند.
زمان مرجع، $T_0=L_0/U_0$ انتخاب می‌شود.
هر مؤلفهٔ حالت، مقیاس متناسب خود را دارد.
دامنهٔ هر تحریک، بر دامنهٔ مرجع همان تحریک تقسیم می‌شود.
بردار $\mathbf{a}^{*}$ نیز تحریک‌های بی‌بعد را دربرمی‌گیرد.
با زمان $\tau=t/T_0$، معادلهٔ عمومی چنین نوشته می‌شود:
\begin{equation}
\frac{\mathrm{d}\mathbf{q}^{*}}{\mathrm{d}\tau}
=\mathbf{F}^{*}\!\left(
\mathbf{q}^{*},\mathbf{a}^{*}(\tau),\boldsymbol{\lambda}^{*},\tau
\right),
\qquad
\mathbf{q}^{*}(0)=\mathbf{q}^{*}_{0}.
\label{eq:lit-forced-dimensionless-ode}
\end{equation}
بردار $\boldsymbol{\lambda}^{*}$، پارامترهای بی‌بعدِ معلوم را نشان می‌دهد.
در این صورت‌بندی، حالت باید برای ادامهٔ حرکت کافی باشد.
در مدل اینرسی‌دار، مکان و سرعت هر دو اهمیت دارند.
اگر مدل حافظه داشته باشد، متغیرهای متناظر نیز لازم هستند.

برای پیش‌بینی زمانی، دو راهبرد اصلی قابل تفکیک هستند.
راهبرد نخست، پاسخ کل بازه را مستقیماً تقریب می‌زند.
در آن، ورودی تابعی و زمان پرس‌وجو وارد شبکه می‌شوند.
راهبرد دوم، یک نگاشت تحول کوتاه‌مدت را می‌آموزد.
سپس، پاسخ کامل با اتصال چند بازه ساخته می‌شود.
کاهش طول بازه، به‌تنهایی تضمین پایداری ایجاد نمی‌کند.

وانگ و پردیکاریس، در سال ۲۰۲۳ راهبرد دوم را بررسی کردند.
آنان، نگاشت حالت اولیه به پاسخ کوتاه‌مدت را یاد گرفتند.
آموزش، با معادلات و بدون زوج‌های برچسب‌دار پاسخ انجام می‌شد.
انتهای هر بازه، حالت اولیهٔ بازهٔ بعد را تأمین می‌کرد.
مثال‌ها، معادلات تحول و دستگاه‌های سخت را نیز شامل می‌شدند.
این مطالعه، مبنایی برای اتصال تکراری عملگرهای محلی فراهم می‌کند.
اما، اعتبار آن برای تحریک تازه باید جدا بررسی شود.
\cite{Wang2023LongTime}

در مسئلهٔ تحریک متغیر، هر پنجره به اطلاعات ورودی خودش نیاز دارد.
برای پنجرهٔ $n$، زمان و تحریک محلی را چنین تعریف می‌کنیم:
\begin{equation}
\begin{aligned}
h_n^{*}&=\tau_{n+1}-\tau_n>0,
&\sigma&=\frac{\tau-\tau_n}{h_n^{*}}\in[0,1],\\
\mathbf{a}^{*}_{n}(\sigma)
&=\mathbf{a}^{*}(\tau_n+h_n^{*}\sigma),
&\mathbf{q}^{*}_{n}(\sigma)
&=\mathbf{q}^{*}(\tau_n+h_n^{*}\sigma).
\end{aligned}
\label{eq:lit-temporal-window-coordinates}
\end{equation}
ورودی شبکه می‌تواند نمونه‌های تحریک، حالت اولیه و پارامترها باشد.
زمان محلی، نقطهٔ پرس‌وجوی پاسخ داخل پنجره را تعیین می‌کند.
اگر طول پنجره تغییر کند، اطلاعات آن نیز لازم است.
اگر معادله به زمان مطلق وابسته باشد، زمان شروع اهمیت دارد.

انتخاب حسگرهای زمانی، به خانوادهٔ تحریک وابسته است.
پالس باریک یا نوسان سریع، ممکن است میان حسگرها پنهان بماند.
افزایش پرس‌وجوهای زمانی، اطلاعات تحریک حذف‌شده را بازنمی‌گرداند.
روش درون‌یابی تحریک نیز باید همراه تعداد حسگرها مشخص شود.
تعداد حسگرها، تعداد فرمان‌های آموزشی و تعداد پنجره‌ها یکسان نیستند.
همچنین، پنجره‌های یک آزمایش نباید بی‌دلیل نمونه‌های مستقل فرض شوند.

اعمال دقیق حالت اولیه، می‌تواند با تبدیل خروجی انجام شود.
برای نمونه، ساختار توضیحی زیر این شرط را اعمال می‌کند:
\begin{equation}
\begin{aligned}
\widehat{\mathbf{q}}^{*}_{n,\theta}(\sigma)
&=\widehat{\mathbf{q}}^{*}_{n,0}
+\sigma\,\mathbf{C}_{\theta}
\left[\mathbf{a}^{*}_{n},\widehat{\mathbf{q}}^{*}_{n,0},
\boldsymbol{\lambda}^{*},h_n^{*},\tau_n\right](\sigma),\\
\widehat{\mathbf{q}}^{*}_{n+1,0}
&=\widehat{\mathbf{q}}^{*}_{n,\theta}(1).
\end{aligned}
\label{eq:lit-temporal-initial-transform}
\end{equation}
در اینجا، $\mathbf{C}_{\theta}$ تصحیح آموخته‌شدهٔ حالت است.
در اجرا، تابع تحریک با نمایش مشخصی وارد مدل می‌شود.
در $\sigma=0$، مقدار اولیه بدون جریمهٔ تقریبی برقرار است.
می‌توان جملهٔ پایهٔ فیزیکی نیز به این ساختار افزود.
این انتخاب، به طراحی مدل و معادلهٔ مسئله وابسته است.

بی‌بعدسازی زمان پنجره، باید در باقیمانده نیز رعایت شود.
برای ساختار بالا، باقیماندهٔ محلی به‌صورت زیر است:
\begin{equation}
\mathbf{R}^{*}_{n,\theta}(\sigma)
=\frac{\partial\widehat{\mathbf{q}}^{*}_{n,\theta}}{\partial\sigma}
-h_n^{*}\mathbf{F}^{*}\!\left(
\widehat{\mathbf{q}}^{*}_{n,\theta}(\sigma),
\mathbf{a}^{*}_{n}(\sigma),\boldsymbol{\lambda}^{*},
\tau_n+h_n^{*}\sigma\right).
\label{eq:lit-temporal-window-residual}
\end{equation}
حذف ضریب $h_n^{*}$، معادلهٔ مورد آموزش را تغییر می‌دهد.
اعمال حالت اولیه نیز، به‌تنهایی باقیمانده را صفر نمی‌کند.
بنابراین، پیوستگی اتصال و رعایت معادله، دو شرط متفاوت هستند.

موضوع دیگری که اهمیت دارد، علیت زمانی است.
پاسخ تا یک زمان، نباید به تغییر تحریک آینده وابسته باشد.
لیو و همکاران، در سال ۲۰۲۴ معماری علّی ارائه کردند.
روش آنان، عملگرهای خطی و پاسخ سازه به زلزله را بررسی کرد.
وزن‌دهی کانولوشنی، وابستگی مجاز زمانی را در ساختار محدود می‌کرد.
این نتیجه، مستقیماً تضمین دستگاه مغناطیسی غیرخطی محسوب نمی‌شود.
\cite{Liu2024Causality}

دریافت کل فرمان، در پیش‌بینی برون‌خط می‌تواند مجاز باشد.
اما، این دسترسی با اعمال ساختاری علیت یکسان نیست.
برای آزمون، می‌توان دو فرمان با گذشتهٔ یکسان ساخت.
سپس، فقط ادامهٔ آیندهٔ یکی از فرمان‌ها تغییر داده می‌شود.
پاسخ‌های پیش از تغییر، باید مطابق علیت یکسان بمانند.
در پیش‌بینی برخط، اطلاعات اندازه‌گیری‌نشدهٔ آینده نیز نباید وارد شود.

در اتصال پنجره‌ها، خطای حالت ورودی می‌تواند منتقل شود.
خطای محلی شبکه نیز به این خطا افزوده می‌شود.
با فرض لیپشیتزبودن نگاشت دقیق، یک کران توضیحی چنین است:
\begin{equation}
e_{n+1}\leq L_{\Phi,n}e_n+\varepsilon_n,
\qquad
e_n=\left\|\widehat{\mathbf{q}}^{*}_{n,0}
-\mathbf{q}^{*}(\tau_n)\right\|_2.
\label{eq:lit-temporal-error-propagation}
\end{equation}
نماد $L_{\Phi,n}$، حساسیت نگاشت تحول به حالت اولیه است.
نماد $\varepsilon_n$، خطای محلی برای حالت ورودی جاری است.
این کران، تحریک و پارامترهای یکسان را در دو حل فرض می‌کند.
کوچکی خطای یک پنجره، کوچکی خطای تمام زنجیره را تضمین نمی‌کند.
همچنین، حالت‌های پیش‌بینی‌شده ممکن است از دامنهٔ آموزش خارج شوند.

ماندل و همکاران، در سال ۲۰۲۶ روش \lr{PITI-DeepONet} را ارائه کردند.
این روش، عملگر مشتق زمانی را با قیود فیزیکی می‌آموزد.
دو خروجی، بازسازی حالت و مشتق زمانی را فراهم می‌کنند.
قید سازگاری، مشتق پیش‌بینی‌شده را به مشتق‌گیری خودکار مرتبط می‌کند.
سپس، انتگرال‌گیر عددی حالت را در زمان پیش می‌برد.
این روش، با پیش‌بینی مستقیم انتهای پنجره تفاوت دارد.
مثال‌های بررسی‌شده، معادلات تحول بودند؛ مدل ذرهٔ مغناطیسی نبودند.
پایش باقیمانده نیز ابزار تشخیص است؛ تضمین عمومی خطا نیست.
\cite{Mandl2026PITI}

برای ذرات مغناطیسی، نوع دادهٔ در دسترس باید روشن باشد.
فن و همکاران، در سال ۲۰۲۶ بستر سیال‌شدهٔ مغناطیسی را بررسی کردند.
آنان، پیش‌بینی مکان ذرات را با داده‌های \lr{CFD-DEM} مطالعه کردند.
در کد عمومی همراه، سرعت، نیرو و تاریخچهٔ مکان ورودی هستند.
این اطلاعات، از دو فرمان مغناطیسی و شرایط اولیه غنی‌تر هستند.
بنابراین، این مسئله با نگاشت تنها فرمان به مسیر یکسان نیست.
ارزیابی با تاریخچهٔ واقعی، با بازپخش کاملاً آزاد متفاوت است.
\cite{Fan2026MagneticBed}

در مسئلهٔ دو منبع، شکل فضایی منابع می‌تواند ثابت باشد.
در این حالت، دو دامنهٔ زمانی، تحریک‌های قابل تغییر هستند.
پارامترهای ذره و جریان حامل نیز باید معلوم یا ورودی باشند.
در مدل بیش‌میرای مقرر، مکان می‌تواند برای ادامهٔ حرکت کافی باشد.
در مدل اینرسی‌دار، بردار سرعت نیز باید منتقل شود.
معادلهٔ نیرو و فرض‌های هر مدل، در فصل سوم مشخص می‌شوند.

جدول~\ref{tab:lit-temporal-operator-studies}، تفاوت نقش مطالعات منتخب را نشان می‌دهد.
این مقایسه، مبنای انتقال ایدهٔ روش‌شناختی میان مسائل است.
از آن، برتری عددی مستقیم میان دستگاه‌های متفاوت نتیجه نمی‌شود.

\begin{table}[htbp]
\centering
\caption{مطالعات منتخب در یادگیری تحول و پاسخ زمانی}
\label{tab:lit-temporal-operator-studies}
\begin{tabular}{p{0.20\textwidth}p{0.31\textwidth}p{0.37\textwidth}}
\hline
مطالعه & ایدهٔ مرتبط & محدودیت انتقال به مسئلهٔ مغناطیسی \\
\hline
وانگ و پردیکاریس، ۲۰۲۳ & اتصال نگاشت‌های تحول کوتاه‌مدت & تعمیم تحریک باید جدا سنجیده شود \\
لیو و همکاران، ۲۰۲۴ & اعمال ساختاری علیت زمانی & تمرکز بر دستگاه‌های خطی \\
ماندل و همکاران، ۲۰۲۶ & یادگیری مشتق و انتگرال‌گیری عددی & تفاوت با پیش‌بینی مستقیم مسیر \\
فن و همکاران، ۲۰۲۶ & پیش‌بینی ذرات بستر مغناطیسی & استفاده از دادهٔ حرکت و نیرو \\
\hline
\end{tabular}
\end{table}

در اعتبارسنجی، دو آزمایش محلی و زنجیره‌ای باید جدا باشند.
آزمایش محلی، می‌تواند هر پنجره را از حالت مرجع آغاز کند.
در آزمایش زنجیره‌ای، حالت پیش‌بینی‌شده به پنجرهٔ بعد منتقل می‌شود.
دومی، انباشت خطا در کاربرد واقعی زنجیره را آشکار می‌کند.
تأمین نیرو یا تاریخچهٔ مرجع نیز باید صریح گزارش شود.
مدت بازپخش، تعداد پنجره‌ها و خانوادهٔ تحریک باید معلوم باشند.
برای مسیر ذره، خطای محوری و عرضی نیز جدا بررسی شوند.
بازپخش مستقل، باید همان فرمان و شرایط اولیه را استفاده کند.
دقت زمانی حل مرجع نیز، با بررسی همگرایی کنترل شود.
این معیارها، پیوند پیش‌بینی مستقیم و تخمین تحریک را روشن می‌کنند.

\قسمت{طراحی رهاسازی و تشکیل آرایش چندذره‌ای}
\label{sec:lit-release-formation}

تشکیل آرایش، به هماهنگی حرکت چند ذره نیاز دارد.
این هماهنگی می‌تواند با کنترل جریان یا طراحی شرایط شروع ایجاد شود.
بااین‌حال، متغیرهای قابل تنظیم و معیار موفقیت روش‌ها متفاوت هستند.
در ادامه، این تفاوت‌ها مبنای مقایسهٔ پژوهش‌های مرتبط قرار می‌گیرند.

\زیرقسمت{هدایت ذرات و طراحی موقعیت نهایی}
\label{sec:lit-particle-terminal-control}

در هدایت ذره، هدف می‌تواند رسیدن به یک موقعیت مشخص باشد.
در رهگیری، علاوه بر مقصد، حرکت میان نقاط نیز اهمیت دارد.
در تشکیل آرایش، چند موقعیت هدف باید هماهنگ تأمین شوند.
بنابراین، مقایسهٔ روش‌ها به تعریف دقیق هدف نیاز دارد.

شنوی و همکاران، تلهٔ استوکس را برای دست‌کاری چندذره‌ای توسعه دادند.
آنان، کنترل هم‌زمان دو ذره را به‌صورت تجربی نشان دادند.
کنترلگر پیش‌بین، دبی درگاه‌های دستگاه را هنگام حرکت تنظیم می‌کرد.
جابه‌جایی مقصد و حرکت روی مسیرهای برنامه‌ریزی‌شده نیز بررسی شد.
در نمونهٔ مونتاژ، دو ذره برای اتصال به یکدیگر نزدیک شدند.
این مطالعه، نمونه‌ای بنیادی از کنترل موقعیت با جریان است.
بااین‌حال، میدان حامل در طول کنترل تغییر می‌کند.
\cite{Shenoy2016StokesTrap}

فانگ و همکاران، مدل حرکت را از داده آموختند.
سپس، مدل آموخته‌شده در کنترلگر نمونه‌برداری‌محور به‌کار رفت.
مثال‌ها، مونتاژ هدفمند، حرکت گروه و عبور از موانع بودند.
این نمایش‌ها، در محیط میکروسیالیِ شبیه‌سازی‌شده انجام شدند.
اهمیت روش، اتصال مدل جانشین به تصمیم‌گیری کنترلی است.
اما، داده‌محوربودن مدل، به‌تنهایی به معنای فیزیک‌آگاه‌بودن آن نیست.
\cite{Fang2023DataDriven}

کیسلایا و همکاران، در سال ۲۰۲۵ جریان قابل‌برنامه‌ریزی را بررسی کردند.
آنان، محفظهٔ هله--شاو را با سه درگاه عمودی ترکیب کردند.
تزریق یا برداشت سیال، خطوط جریان را تغییر می‌داد.
مسیرها از پیش بهینه می‌شدند و بازخورد، انحراف تجربی را اصلاح می‌کرد.
هزینهٔ طراحی، اندازهٔ دبی‌ها و تغییرات زمانی آن‌ها را دربرمی‌گرفت.
کاربردها، به‌دام‌اندازی، جداسازی و نزدیک‌کردن ذرات بودند.
در اینجا، درگاه کنترل جریان با ایستگاه رهاسازی ذره متفاوت است.
\cite{Kislaya2025Programmable}

این مطالعات، چند معیار متفاوت برای عملکرد مطرح می‌کنند.
خطای مقصد، زمان انجام حرکت و میزان تحریک، معیارهای مکمل هستند.
در کنترل بازخوردی، مشاهدهٔ انحراف می‌تواند فرمان بعدی را تغییر دهد.
برنامهٔ رهاسازی ازپیش‌محاسبه‌شده، در اجرای بدون بازخورد چنین اصلاحی ندارد.
بنابراین، حساسیت به خطای شرایط رهاسازی نیز باید بررسی شود.
مقایسهٔ سرعت روش‌ها، بدون احتساب مشاهده و محاسبه، کامل نیست.

\زیرقسمت{طراحی مکان، سرعت و زمان رهاسازی}
\label{sec:lit-release-design-variables}

در طراحی رهاسازی، انتخاب‌ها پیش از ورود ذره به جریان انجام می‌شوند.
مکان شروع و زمان ورود، دو متغیر مهم هستند.
در مدل اینرسی‌دار، سرعت اولیه نیز متغیر طراحی است.
در مدل بیش‌میرا، معمولاً چنین آزادی مستقلی وجود ندارد.
بازهٔ مجاز هر متغیر، باید با مدل و امکانات دستگاه سازگار باشد.

راه دیگر، تغییر ساختار خود دستگاه برای دستیابی به مسیر مطلوب است.
لی و همکاران، در سال ۲۰۲۶ طراحی شبکهٔ کانال را بررسی کردند.
چارچوب آنان، پیش‌بینی محلی حرکت را میان ماژول‌های کانال ترکیب می‌کرد.
طراحی معکوس، ساختارها و تنظیمات متناسب با رفتار مطلوب را جست‌وجو می‌کرد.
این روش، اهمیت اطلاعات موقعیت، سرعت و زمان عبور را نشان می‌دهد.
اما، تغییر هندسه با تغییر رهاسازی در هندسهٔ ثابت یکسان نیست.
\cite{Li2026ParticleArchitecture}

برای شرایط شروع، مقیاس‌های رابطهٔ \ref{eq:lit-reference-scales} به‌کار می‌روند.
سرعت ذره نیز مطابق بخش \ref{sec:lit-lagrangian-particle} بی‌بعد می‌شود.
برای ذرهٔ $i$، قطر بی‌بعد با $d_i^*=d_i/L_0$ تعریف می‌شود.
زمان بی‌بعد را با $\tau=t/T_0$ نشان می‌دهیم.
این تعریف، هندسه یا معادلهٔ حرکت را تغییر نمی‌دهد.

در حالت ایستگاه ثابت، مکان شروع از مجموعه‌ای محدود انتخاب می‌شود.
در حالت مکان آزاد، موقعیت روی مرز نیز قابل تنظیم است.
یک ایستگاه، در صورت اجازهٔ قیود، می‌تواند چند رهاسازی انجام دهد.
پس، تعداد ایستگاه‌ها با تعداد ذرات برابر نیست.
تعداد رهاسازی مجاز و حداقل فاصلهٔ زمانی نیز قیود جدا هستند.
در مقایسهٔ این حالت‌ها، حدود سرعت و بودجهٔ رهاسازی باید مشخص باشند.

زمان مشاهدهٔ مشترک را با $T^{*}$ نشان می‌دهیم.
زمان رهاسازی ذرهٔ $i$ نیز $\tau_i$ است.
سن پرواز، مدت حرکت از رهاسازی تا زمان مشاهده است.
برای حرکت مستقل در میدان پایا، نگاشت توضیحی زیر را درنظر می‌گیریم:
\begin{equation}
\begin{aligned}
\Delta\tau_i&=T^{*}-\tau_i\geq0,\\
\mathbf{r}_i^{*}(T^{*})
&=\Phi^{*}_{\Delta\tau_i}
\left(\mathbf{r}_{0,i}^{*},\mathbf{v}_{0,i}^{*};
\boldsymbol{\lambda}^{*}\right).
\end{aligned}
\label{eq:lit-release-common-clock}
\end{equation}
بردار $\boldsymbol{\lambda}^{*}$، پارامترهای معلوم مدل را دربرمی‌گیرد.
در مدل بیش‌میرا، ورودی سرعت اولیهٔ مستقل حذف می‌شود.
برای میدان ناپایا، زمان مطلق یا تاریخچهٔ تحریک نیز لازم است.
مقایسهٔ انتهای مسیرها با مدت برابر، لزوماً آرایش هم‌زمان نیست.
ذراتِ رهاشده در زمان‌های متفاوت، در یک لحظه سن متفاوت دارند.

هم‌زمانی، قیدی فراتر از دقت مقصد هر ذره ایجاد می‌کند.
باید همهٔ ذرات، در همان لحظه به اهداف متناظر نزدیک باشند.
وقتی مسیر زودتر خاتمه یابد، رسیدن به هدف پذیرفته نیست.
تماس با دیواره یا خروج، باید مطابق مدل رویداد تفسیر شود.
ادامه‌دادن مصنوعی مسیر خاتمه‌یافته، گواه تشکیل آرایش نیست.

همهٔ نقاط هدف، الزاماً از هر ایستگاه دسترس‌پذیر نیستند.
زمان حرکت و حدود سرعت اولیه، مجموعهٔ قابل‌دستیابی را تغییر می‌دهند.
اگر ایستگاه‌های قبلی حفظ شوند، گزینه‌های طراحی کمتر نمی‌شوند.
اما، خطای عددیِ بهینه‌ساز الزاماً به‌صورت یکنواخت کاهش نمی‌یابد.
بودجهٔ محاسباتی و شروع‌های اولیه نیز باید مشخص باشند.
این تفکیک، برای تحلیل اثر تعداد ایستگاه‌ها اهمیت دارد.

\زیرقسمت{تشکیل آرایش و قیود چندذره‌ای}
\label{sec:lit-multiparticle-formation-constraints}

شکل هدف را می‌توان با نقاطی برای مراکز ذرات نمایش داد.
در این نمایش، تعداد نقاط و اندازهٔ ذرات اهمیت دارند.
فاصلهٔ بسیار کم اهداف، می‌تواند تشکیل دقیق آرایش را ناممکن کند.
برای ذرات هم‌نقش، تخصیص ذره به هدف نیز قابل انتخاب است.
یک معیار توضیحیِ خطای نهایی چنین تعریف می‌شود:
\begin{equation}
\min_{\Theta\in\mathcal{A},\,\pi\in\mathfrak{S}_N}
E_T(\Theta,\pi),
\qquad
E_T=\frac{1}{N}\sum_{i=1}^{N}
\left\|\mathbf{r}_i^{*}(T^{*};\Theta)
-\mathbf{z}^{*}_{\pi(i)}\right\|_2^2.
\label{eq:lit-formation-assignment-objective}
\end{equation}
نماد $\Theta$، متغیرهای طراحی و $\mathcal{A}$، مجموعهٔ انتخاب‌های مجاز هستند.
نقاط $\mathbf{z}^{*}_j$، مراکز هدف و $N$، تعداد ذرات هستند.
جایگشت $\pi$، تخصیص یک‌به‌یک ذرات به نقاط هدف را مشخص می‌کند.
اگر ذرات نقش متفاوت داشته باشند، تخصیص نیز محدود می‌شود.
این تعریف، جانشین جزئیات تابع هزینهٔ فصل سوم نیست.

اربن و همکاران، در سال ۲۰۲۴ دست‌کاری اپتوسیالی را بررسی کردند.
جریان‌های موضعیِ ناشی از لیزر، موقعیت نسبی ذرات را تنظیم می‌کردند.
آنان، کنترل مستقل ۱۵ ذره و ۳۰ درجهٔ آزادی را گزارش کردند.
در روش، تخصیص ذرات به مقصدها نیز بهینه می‌شد.
تولید متوالی و تسهیم زمانی جریان‌ها، نقش مهمی در عملکرد داشت.
این آرایش‌ها، با تحریک مداوم و بازخورد ساخته یا هدایت می‌شدند.
بنابراین، همان مسئلهٔ رهاسازی مستقل در میدان ثابت نیستند.
همچنین، درجهٔ آزادی را نباید با تعداد ذرات اشتباه گرفت.
\cite{Erben2024Optofluidic}

مطالعهٔ شوم و همکاران، نقش برهم‌کنش هیدرودینامیکی را روشن می‌کند.
آنان، کنترل یک ذرهٔ فعال و ذرات منفعل را بررسی کردند.
اثبات کنترل‌پذیری، برای یک یا دو ذرهٔ منفعل ارائه شد.
فرض‌ها، شامل جریان استوکس دوردست و جدایی کافی ذرات بودند.
این نتیجه، به مدل و مجموعهٔ کنترل‌های همان مطالعه وابسته است.
نمی‌توان آن را مستقیماً به آرایش متراکمِ صدذره‌ای تعمیم داد.
\cite{Shum2025HydrodynamicControl}

الهاشم و همکاران، بهینه‌سازی جریان را با مشتق‌گیری خودکار بررسی کردند.
روش، مشتق هدف نسبت به متغیرهای طراحی را محاسبه می‌کرد.
مثال‌ها، حرکت مادهٔ فعال، پخش و اختلاط را شامل می‌شدند.
این مطالعه، ارزش محاسبهٔ گرادیان در طراحی پُربعد را نشان می‌دهد.
بااین‌حال، مشتق‌پذیری حل‌گر به معنای آموزش \lr{PINN} نیست.
همچنین، بهینه‌سازی عددی به‌تنهایی اثبات بهترین جواب سراسری نیست.
\cite{Alhashim2025DifferentiableFlow}

در مسئلهٔ چندذره‌ای، خطای نهایی باید همراه فاصلهٔ ذرات بررسی شود.
تقاطع هندسی مسیرها، الزاماً برخورد هم‌زمان را نشان نمی‌دهد.
فاصله باید در زمان یکسان و حضور مشترک سنجیده شود.
برای مسیرهای معتبر تا $T^{*}$، تعریف توضیحی زیر مناسب است:
\begin{equation}
\begin{aligned}
\mathcal{I}_{ij}
&=[\max(\tau_i,\tau_j),T^{*}],\\
D_{ij}^{*}
&=\min_{\tau\in\mathcal{I}_{ij}}
\left\|\mathbf{r}^{*}_{i}(\tau)-\mathbf{r}^{*}_{j}(\tau)\right\|_2,\\
D_{ij}^{*}
&\geq\frac{d_i^{*}+d_j^{*}}{2}+\delta^{*},
\qquad i\neq j.
\end{aligned}
\label{eq:lit-formation-active-interval-separation}
\end{equation}
در اینجا، $\delta^{*}\geq0$ حاشیهٔ جدایی بی‌بعد است.
این رابطه، برای ذرات کروی یا دایره‌ایِ مدل‌شده تفسیر هندسی دارد.
پیش از رهاسازی، ذره هنوز عضو مجموعهٔ ذرات حاضر نیست.
کنترل فاصله فقط در زمان نهایی، نزدیک‌شدن زودتر را آشکار نمی‌کند.
نمونه‌برداری زمانیِ محدود نیز ممکن است کمینهٔ واقعی را نبیند.
پس، کنترل متراکم‌تر و پالایش کمینه، شواهد عددی مهمی فراهم می‌کنند.

جریمهٔ فاصله، معادل حل کامل برخورد نیست.
اثر تماس، روانکاری و نوسان براونی، به مدل جدا نیاز دارند.
فرض عدم تأثیر بر میدان، باید از حذف برهم‌کنش جدا شود.
کنارگذاشتن این اثرها، باید هنگام تفسیر آرایش متراکم روشن باشد.
نتیجهٔ دوبعدی نیز، خودبه‌خود تضمین آرایش سه‌بعدی نیست.

تشکیل شکل در یک لحظه، با حفظ آن در زمان متفاوت است.
برای ادعای پایداری، بازهٔ پس از تشکیل نیز باید بررسی شود.
اگر آرایش جابه‌جا شود، مرجع مکانی نیز باید مشخص باشد.
اتصال چسبندهٔ ذرات نیز، سازوکاری متفاوت برای حفظ ساختار است.
بنابراین، مونتاژ متصل، تثبیت بازخوردی و آرایش گذرا نباید یکسان گزارش شوند.

جدول~\ref{tab:lit-formation-design-comparison}، تفاوت متغیرهای طراحی را خلاصه می‌کند.
این جدول، رتبه‌بندی دقت عددی میان دستگاه‌های متفاوت نیست.

\begin{table}[htbp]
\centering
\caption{مقایسهٔ متغیرهای طراحی در مطالعات مرتبط با هدایت و آرایش}
\label{tab:lit-formation-design-comparison}
\begin{tabular}{p{0.19\textwidth}p{0.24\textwidth}p{0.21\textwidth}p{0.24\textwidth}}
\hline
مطالعه & متغیر قابل تنظیم & نوع هدف & تفاوت با طراحی رهاسازی \\
\hline
شنوی، ۲۰۱۶ & دبی درگاه‌ها با کنترل پیش‌بین & هدایت و اتصال ذرات & تغییر جریان هنگام حرکت \\
فانگ، ۲۰۲۳ & فرمان‌های کنترل هیدرودینامیکی & مسیر و مونتاژ هدفمند & بازخورد با مدل آموخته‌شده \\
کیسلایا، ۲۰۲۵ & دبی سه درگاه و اصلاح بازخوردی & هدایت، جداسازی و نزدیکی & جریان حامل قابل تنظیم \\
لی، ۲۰۲۶ & ماژول‌ها، هندسه و تنظیمات دستگاه & رفتار مکانی و زمانی & طراحی ساختار کانال \\
اربن، ۲۰۲۴ & تحریک‌های اپتوسیالی موضعی & موقعیت و آرایش متحرک & تحریک مداوم و تخصیص اهداف \\
شوم، ۲۰۲۵ & حرکت ذرهٔ فعال & کنترل ذرات منفعل & اتکا به برهم‌کنش دوردست \\
الهاشم، ۲۰۲۵ & پارامترهای طراحی در حل‌گر مشتق‌پذیر & اهداف مختلف جریان & روش عمومی بهینه‌سازی \\
\hline
\end{tabular}
\end{table}

برای ارزیابی طراحی رهاسازی، چند گزارش باید کنار هم ارائه شوند.
تعداد ذرات، ایستگاه‌ها و رهاسازی‌های مجاز باید جدا ثبت شوند.
خطای میانگین و بیشینهٔ خطای مقصد، رفتارهای مکمل را نشان می‌دهند.
کمینهٔ فاصله و اعتبار مسیرها نیز باید گزارش شوند.
بازپخش مستقل برنامه، می‌تواند ناسازگاری مدل جانشین را آشکار کند.
بودجهٔ بهینه‌سازی و نرخ موفقیت تکرارها نیز اهمیت دارند.
افزایش تراکم یا پیچیدگی شکل، باید همراه همین اطلاعات تفسیر شود.
این معیارها، جایگاه طراحی رهاسازی را میان روش‌های هدایت روشن می‌کنند.
کمبود نمونهٔ کاملاً مشابه، به‌تنهایی اثبات نوآوری پژوهش نیست.

\قسمت{بازسازی جریان و تحریک از حرکت ذرات}
\label{sec:lit-trajectory-inverse}

حرکت ذرات، اطلاعاتی غیرمستقیم از جریان و نیروهای مؤثر فراهم می‌کند.
استخراج این اطلاعات، به مدل حرکت و نوع مشاهدات وابسته است.
در بازسازی، مجهول می‌تواند میدان جریان یا تحریک واردشده باشد.
این دو مسئله، داده‌ها و محدودیت‌های یکسانی ندارند.
ابتدا، بازسازی جریان از مسیرهای زمان‌دار بررسی می‌شود.

\زیرقسمت{بازسازی جریان از مسیرهای زمان‌دار}
\label{sec:lit-timed-flow-inverse}

\textbf{دادهٔ لاگرانژی و مجهول اویلری}

در توصیف لاگرانژی، حرکت یک ذره در زمان دنبال می‌شود.
در توصیف اویلری، میدان در مختصات مکانی و زمانی بیان می‌شود.
بازسازی جریان، پیوندی میان این دو نوع اطلاعات ایجاد می‌کند.
دادهٔ مسیر زمان‌دار، شامل مکان ذره و زمان ثبت آن است.
بنابراین، شکل هندسی مسیر به‌تنهایی تمام دادهٔ مسئله نیست.
یک منحنی یکسان، می‌تواند با سرعت‌های متفاوت پیموده شود.
زمان ثبت، این تفاوت را در مشاهدات قابل تشخیص می‌کند.

کمیت‌ها، با مقیاس‌های رابطهٔ \ref{eq:lit-reference-scales} بی‌بعد می‌شوند.
سرعت ذره نیز بر مقیاس ثابت $U_0$ تقسیم می‌شود.
پس از این تعریف، ستاره‌ها برای سادگی حذف می‌شوند.
تمام روابط بعدی این زیربخش، با کمیت‌های بی‌بعد نوشته می‌شوند.
نماد $\mathbf r_i(t_{ij})$ مکان ذرهٔ $i$ در زمان $t_{ij}$ است.
تعداد زمان‌های ثبت، می‌تواند برای مسیرهای مختلف متفاوت باشد.
تعداد ذرات، تعداد نقاط زمانی و مدت مشاهده، معیارهای جداگانه‌اند.
گزارش تعداد ذرات، بدون مشخصات ثبت زمانی، توصیف کاملی نیست.

\textbf{ردیاب و ذرهٔ اینرسی‌دار}

نوع ذره، رابطهٔ میان حرکت مشاهده‌شده و جریان را تعیین می‌کند.
برای ردیاب ایدئال، سرعت ذره برابر سرعت موضعی سیال است.
برای ذرهٔ اینرسی‌دار، اختلاف این دو سرعت باید مدل شود.
مدل کامل‌تر حرکت، می‌تواند چندین نیروی هیدرودینامیکی را دربرگیرد.
پس انتخاب مدل، بخشی از تعریف مسئلهٔ معکوس است.
\cite{MaxeyRiley1983}

برای روشن‌شدن تفاوت، تنها درگ خطی استوکس را در نظر می‌گیریم.
زمان پاسخ ذره ثابت است و نیروی دیگری اعمال نمی‌شود.
در این حالت، دو مدل مشاهده به‌صورت زیر نوشته می‌شوند.
\begin{equation}
\begin{aligned}
\text{ردیاب:}\qquad
&\frac{\mathrm d\mathbf r}{\mathrm dt}
=\mathbf u(\mathbf r(t),t),\\
\text{ذرهٔ اینرسی‌دار:}\qquad
&\frac{\mathrm d\mathbf r}{\mathrm dt}=\mathbf v_p,
\qquad
\mathrm{St}\frac{\mathrm d\mathbf v_p}{\mathrm dt}
=\mathbf u(\mathbf r(t),t)-\mathbf v_p,\\
&\mathrm{St}=\frac{\tau_p}{T_0}>0.
\end{aligned}
\label{eq:lit-timed-observation-models}
\end{equation}
در این رابطه، $\tau_p$ زمان پاسخ بُعددار ذره است.
عدد استوکس، نسبت زمان پاسخ به مقیاس زمان مرجع است.
پس مشتق مسیر اینرسی‌دار، مستقیماً سرعت سیال را نشان نمی‌دهد.
در مدل ساده، شتاب ذره نیز برای تعیین سرعت سیال لازم است.
مشتق دوم عددی مکان، خطای اندازه‌گیری را تشدید می‌کند.
هموارسازی می‌تواند نویز را کاهش دهد، اما جزئیات حرکت را نیز می‌کاهد.
مشتق‌گیری خودکار نیز خطای برازش دادهٔ نویزی را حذف نمی‌کند.
\cite{ZhouGrauer2026Joint}

\textbf{نوع مشاهده در پژوهش‌های بازسازی}

رئیسی و همکاران، چارچوب مکانیک سیالات پنهان را معرفی کردند.
در این روش، تصاویر غلظت یک اسکالر غیرفعال استفاده می‌شوند.
معادلهٔ انتقال اسکالر، همراه ناویر–استوکس، تخمین سرعت و فشار را مقید می‌کند.
پس دادهٔ ورودی، مجموعه‌ای از مسیرهای منفرد ذرات نیست.
این مطالعه، اهمیت تعریف دقیق مشاهده را نشان می‌دهد.
قابلیت استنتاج میدان، به اطلاعات موجود در تصاویر نیز وابسته است.
بنابراین، حذف برخی شرایط مرزی، مستقل از کفایت داده نیست.
این چارچوب، مبنایی برای بازسازی فیزیک‌آگاه از اطلاعات غیرمستقیم فراهم کرد.
\cite{Raissi2020HFM}

اشتاینفورت و همکاران، میدان حلقه‌های گردابه را از دادهٔ رهگیری بازسازی کردند.
آنان، سرعت‌های حاصل از رهگیری سه‌بعدی ذرات را به شبکه ارائه کردند.
قیود ناویر–استوکس، بازسازی میدان میان مشاهدات پراکنده را هدایت می‌کردند.
حالت‌های دادهٔ تجمیع‌شده و تک‌برداشت نیز بررسی شدند.
این ورودی، با دریافت مستقیم موقعیت‌های خام یکسان نیست.
در نتیجه، نوع پردازش مسیر باید در مقایسه ثبت شود.
\cite{Steinfurth2024}

\textbf{یادگیری مسیر و میدان از دادهٔ زمان‌دار}

تاکور و رئیسی، چارچوب \lr{ELPINN} را در پیش‌چاپی معرفی کردند.
مشاهدهٔ اولیه، موقعیت زمان‌دار ردیاب‌های بدون اینرسی است.
ابتدا، نگاشت مسیر با استفاده از موقعیت‌ها آموزش می‌بیند.
سپس، مشتق زمانی این نگاشت، برای یادگیری سرعت استفاده می‌شود.
در نمونه‌های دوبعدی، فشار با باقیماندهٔ مومنتوم تخمین زده می‌شود.
این مراحل، با ثابت‌کردن مدل‌های قبلی انجام می‌شوند.
بنابراین، ساختار ارائه‌شده، آموزش مشترک تمام مؤلفه‌ها نیست.
نویسندگان، اثر تفکیک زمانی و تعداد داده‌ها را بررسی کردند.
در محدودهٔ آزمون‌ها، کاهش فاصلهٔ زمانی تأثیر محسوسی داشت.
بااین‌حال، این نتیجه، بی‌اهمیت‌بودن تعداد مسیرها را اثبات نمی‌کند.
\cite{Thakur2025ELPINN}

وان و همکاران، مدل \lr{TrajectoryFlowNet} را برای یادگیری مشترک پیشنهاد کردند.
این مدل، دو بخش برای مسیر و میدان جریان دارد.
سازگاری مشتق زمانی مسیر با سرعت، این بخش‌ها را پیوند می‌دهد.
فرض حرکت، ردیاب غیرفعال با اینرسی ناچیز است.
نسخهٔ بررسی‌شده، از خطای موقعیت و خطای سرعت استفاده می‌کند.
پس نتایج آن، تنها بر دادهٔ موقعیت خام متکی نیستند.
در آزمون حفره، مجموعهٔ پایه شامل ۲۰۰ مسیر آموزشی بود.
در مکمل، سهم داده تا ده درصد کاهش یافت.
کاهش داده، خطای مسیر را بیش از میدان افزایش داد.
این آزمون، حداقل عمومی تعداد ذرات را تعیین نمی‌کند.
\cite{Wan2026TrajectoryFlowNet}

این دو رویکرد، نقش نگاشت مسیر را برجسته می‌کنند.
بااین‌حال، مشاهده‌ها و ترتیب آموزش آن‌ها یکسان نیست.
همچنین، فرض ردیاب با مدل ذرهٔ اینرسی‌دار تفاوت دارد.
مقایسهٔ تعداد ذرات، باید پس از روشن‌کردن این تفاوت‌ها انجام شود.
در غیر این صورت، دشواری دو مسئله درست سنجیده نمی‌شود.

\textbf{بازسازی با لحاظ اینرسی و روش‌های غیرشبکه‌ای}

ژو و گراور، تخمین مشترک جریان و ویژگی‌های ذرات را بررسی کردند.
چارچوب \lr{NIPA}، میدان اویلری را به مدل‌های لاگرانژی ذرات متصل می‌کند.
مشاهدات، موقعیت‌های زمان‌دار هستند و عدم‌قطعیت مکان‌یابی نیز لحاظ می‌شود.
مدل مسیر، سازگاری سینماتیکی را در ساختار خود اعمال می‌کند.
قیود جریان و دینامیک ذره، در تخمین مشترک مشارکت دارند.
مطالعه، ردیاب‌ها و ذرات اینرسی‌دار را در چند رژیم بررسی می‌کند.
در برخی آزمون‌ها، اندازه یا چگالی ذره نیز مجهول است.
حساسیت بازسازی به تراکم، نویز و اینرسی گزارش شده است.
این نتایج عددی، اثبات یکتایی عمومی مسئله نیستند.
\cite{ZhouGrauer2026Joint}

بازسازی لاگرانژی، به استفاده از شبکهٔ عصبی محدود نیست.
ناکائو و همکاران، روش تغییراتی چهاربعدی را برای همرفت حرارتی به‌کار گرفتند.
حل مستقیم و حل الحاقی، برای برازش داده‌ها تکرار می‌شوند.
تابع هدف، اختلاف موقعیت و سرعت ذرات را دربرمی‌گیرد.
میدان حرارتی و عدد رایلی نیز در مطالعه تخمین زده می‌شوند.
مدل دوبعدی بسیار لزج، جمله‌های اینرسی سیال را حذف می‌کند.
بنابراین، نتایج آن به هر رژیم ناویر–استوکس تعمیم نمی‌یابند.
این نمونه، یک مقایسهٔ روش‌شناختی با بازسازی شبکه‌ای فراهم می‌کند.
\cite{Nakao2025LagrangianDA}

در مجموع، قید فیزیکی می‌تواند در ساختار یا تابع هدف وارد شود.
حل‌گر عددی نیز می‌تواند معادلات گسسته را در بازسازی اعمال کند.
این انتخاب‌ها، هزینهٔ محاسباتی و خطاهای متفاوتی ایجاد می‌کنند.
هیچ‌کدام، به‌تنهایی کفایت مشاهدات یا همگرایی تخمین را تضمین نمی‌کنند.

\textbf{صورت‌بندی انتگرالی و استفاده از موقعیت}

برای مدل درگ خطی، می‌توان رابطه‌ای مستقیم برای موقعیت نوشت.
شرایط اولیه، $\mathbf r_0$ و $\mathbf v_{p,0}$ در زمان $t_0$ هستند.
حل معادلهٔ سرعت و انتگرال‌گیری زمانی، رابطهٔ زیر را می‌دهد.
\begin{equation}
\begin{aligned}
\mathbf r(t)={}&\mathbf r_0
+\mathrm{St}\left[1-\exp\!\left(-\frac{t-t_0}{\mathrm{St}}\right)\right]
\mathbf v_{p,0}\\
&+\int_{t_0}^{t}
\left[1-\exp\!\left(-\frac{t-s}{\mathrm{St}}\right)\right]
\mathbf u(\mathbf r(s),s)\,\mathrm ds.
\end{aligned}
\label{eq:lit-timed-integral-position}
\end{equation}
این رابطه، تنها برای فرض‌های معادلهٔ \ref{eq:lit-timed-observation-models} معتبر است.
برای درگ غیرخطی، همین هستهٔ نمایی عموماً معتبر نیست.
صورت انتگرالی، مقایسه با موقعیت را بدون مشتق دوم عددی ممکن می‌کند.
بااین‌حال، نویز مکان و خطای مدل همچنان باقی می‌مانند.
درگ، زمان پاسخ و شرایط اولیه نیز باید درست مشخص شوند.
درون‌یابی زمانی مشاهدات، دادهٔ مستقل تازه‌ای ایجاد نمی‌کند.
همچنین، حذف مشتق دوم، معادل اثبات پایداری کامل تخمین نیست.

هنگام برازش انتگرالی، می‌توان میدان را روی مسیر مشاهده‌شده ارزیابی کرد.
در این حالت، محل نمونه‌برداری میدان از داده گرفته می‌شود.
در بازپخش، مکان ذره از حل معادلهٔ حرکت به‌دست می‌آید.
پس محل ارزیابی میدان نیز همراه پاسخ تخمینی تغییر می‌کند.
این تفاوت، بازپخش را به آزمونی ضروری برای سازگاری مسیر تبدیل می‌کند.
بااین‌حال، بازپخش موفق، به‌تنهایی دقت سراسر میدان را اثبات نمی‌کند.

\textbf{دامنهٔ مجهول و تفسیر دادهٔ کم}

تخمین مستقیم میدان، با تخمین ورودی یک خانوادهٔ میدان متفاوت است.
در حالت نخست، شبکه می‌تواند متغیرهای جریان را مستقیماً تقریب بزند.
در حالت دوم، ساختار پاسخ از پیش محدود شده است.
برای نمونه، هندسه و خانوادهٔ پروفایل ورودی می‌توانند معلوم باشند.
آنگاه، مجهول به ضرایب نمایش پروفایل و پارامترهای مجاز محدود می‌شود.
یک عملگر آموزش‌دیده، نگاشت این ورودی‌ها به میدان را تقریب می‌زند.
این نقش، با مفهوم یادگیری عملگر سازگار است.
\cite{lu2021deeponet}

برای نمایش این تفاوت، خانوادهٔ میدان را چنین تعریف می‌کنیم.
\begin{equation}
(\mathbf u_{\boldsymbol\alpha},p_{\boldsymbol\alpha})
=\mathcal G(\boldsymbol\alpha),
\qquad
\boldsymbol\alpha\in\mathcal A\subset\mathbb R^{m}.
\label{eq:lit-timed-restricted-family}
\end{equation}
بردار $\boldsymbol\alpha$، مجهول‌های بی‌بعد خانوادهٔ انتخاب‌شده را نمایش می‌دهد.
این بردار، تمام درجات آزادی یک میدان دلخواه را ندارد.
بنابراین، تعداد مسیرهای لازم می‌تواند به محدودیت خانواده وابسته باشد.
بازسازی با دو مسیر، در چنین چارچوبی قابل بررسی است.
اما موفقیت آن، بازیابی هر میدان دلخواه را اثبات نمی‌کند.
یک مسیر طولانی نیز الزاماً پوشش مکانی مناسبی فراهم نمی‌کند.
مسیرهای نزدیک، ممکن است اطلاعات مشابهی از خانوادهٔ پاسخ منتقل کنند.
حساسیت مسیر به پارامترهای ورودی، از تعداد مسیرها نیز اهمیت می‌گیرد.
پس دادهٔ کم، باید با مشخصات مسئله و پیش‌دانسته‌ها توصیف شود.

\textbf{اعتبارسنجی و جایگاه پژوهش حاضر}

ارزیابی بازسازی، به چند آزمون مکمل نیاز دارد.
نخست، اختلاف پاسخ با موقعیت‌های مشاهده‌شده سنجیده می‌شود.
سپس، ذره از شرایط اولیه در میدان تخمینی بازپخش می‌شود.
این بازپخش، باید با انتگرال‌گیر مستقل نیز کنترل شود.
حل‌گر مستقل، به‌تنهایی استقلال میدان حامل را تضمین نمی‌کند.
بنابراین، میدان باید جداگانه با مرجع معتبر مقایسه شود.
نقاط ارزیابی خارج مسیرها، برای این مقایسه اهمیت دارند.
اگر مجهول پروفایل ورودی باشد، خطای همان پروفایل نیز گزارش می‌شود.
خطای مسیر، خطای ورودی و خطای میدان، معیارهای قابل جایگزینی نیستند.

فشار نیز، معمولاً از معادلات جریان و قیود موجود استنباط می‌شود.
اگر مرجع فشار اعمال نشده باشد، ابهام یک ثابت مکانی باقی می‌ماند.
در جریان ناپایا، این ابهام می‌تواند تابعی از زمان باشد.
اگر شرط مرزی فشار را تعیین کند، همان شرط باید حفظ شود.
پس نحوهٔ مقایسهٔ فشار، به صورت‌بندی واقعی مسئله وابسته است.

جدول \ref{tab:lit-timed-flow-comparison} تفاوت داده‌ها و مجهول‌های پژوهش‌ها را خلاصه می‌کند.
ردیف‌ها، آزمون مشترکی برای رتبه‌بندی دقت روش‌ها تشکیل نمی‌دهند.
تفاوت هندسه، رژیم جریان و پیش‌دانسته‌ها باید در تفسیر حفظ شود.

\begin{table}[htbp]
\centering
\small
\caption{مقایسهٔ مشاهدات و مجهول‌ها در بازسازی جریان از داده‌های غیرمستقیم}
\label{tab:lit-timed-flow-comparison}
\setlength{\tabcolsep}{5pt}
\begin{tabular}{p{0.19\textwidth}p{0.24\textwidth}p{0.25\textwidth}p{0.22\textwidth}}
\hline
پژوهش & دادهٔ اصلی & مجهول یا پاسخ & نکتهٔ مقایسه \\
\hline
رئیسی و همکاران \cite{Raissi2020HFM}
& تصاویر اسکالر غیرفعال
& سرعت و فشار
& مشاهدهٔ اسکالر، بدون مسیر منفرد \\
اشتاینفورت و همکاران \cite{Steinfurth2024}
& سرعت حاصل از رهگیری
& میدان سه‌بعدی و فشار
& ورودی پردازش‌شدهٔ رهگیری \\
تاکور و رئیسی \cite{Thakur2025ELPINN}
& موقعیت زمان‌دار ردیاب
& مسیر، سرعت و فشار دوبعدی
& آموزش مرحله‌ای؛ پیش‌چاپ \\
وان و همکاران \cite{Wan2026TrajectoryFlowNet}
& موقعیت و سرعت ردیاب
& مسیر و میدان جریان
& آموزش مشترک دو بخش \\
ژو و گراور \cite{ZhouGrauer2026Joint}
& موقعیت‌های زمان‌دار
& جریان، مسیر و ویژگی‌های ذره
& آزمون ردیاب و ذرهٔ اینرسی‌دار \\
ناکائو و همکاران \cite{Nakao2025LagrangianDA}
& موقعیت و سرعت ذرات
& جریان، دما و عدد رایلی
& روش تغییراتی، بدون شبکه \\
\hline
\end{tabular}
\end{table}

در پژوهش حاضر، تمرکز بر خانواده‌ای مشخص از جریان کانال است.
پروفایل ورودی، از حرکت تعداد محدودی ذره تخمین زده می‌شود.
هندسه، مدل ذره و اطلاعات اولیه، دامنهٔ این تخمین را محدود می‌کنند.
بنابراین، سهم پژوهش با همین فرض‌ها و آزمون‌های مستقل سنجیده می‌شود.
جزئیات صورت‌بندی و نتایج عددی، در فصل‌های سوم و چهارم ارائه می‌شوند.
در ادامه، اثر حذف زمان اندازه‌گیری‌شده از مسیر بررسی خواهد شد.

\زیرقسمت{بازسازی از هندسهٔ مسیر بدون زمان اندازه‌گیری‌شده}
\label{sec:lit-untimed-flow-inverse}

\textbf{تفاوت مسیر زمان‌دار و منحنی مشاهده‌شده}

در برخی مشاهدات، شکل مسیر معلوم است، اما زمان ثبت نیست.
در این حالت، داده یک منحنی فضایی را مشخص می‌کند.
فرض می‌کنیم نقطهٔ شروع و جهت پیمایش نیز معلوم هستند.
ترتیب نقاط، این داده را از مجموعهٔ نقاط نامرتب متمایز می‌کند.
بااین‌حال، فاصلهٔ زمانی میان نقاط، دیگر مشاهدهٔ مسئله نیست.
بنابراین، نمی‌توان سرعت را مستقیماً از اختلاف مکان‌ها محاسبه کرد.
روش‌های زمان‌دار، این فاصله‌های زمانی را در مدل مشاهده استفاده می‌کنند.
پس حذف زمان، صرفاً تغییر قالب ذخیره‌سازی داده نیست.
\cite{ZhouGrauer2026Joint}

تحلیل هندسی منحنی‌ها، تمایز شکل و پارامتردهی را روشن می‌کند.
سریواستاوا و همکاران، این تمایز را در تحلیل منحنی‌های کشسان بررسی کردند.
در چارچوب آنان، بازپارامتردهی می‌تواند بدون تغییر شکل منحنی انجام شود.
نمایش ریشهٔ دوم سرعت، مقایسهٔ کشسان منحنی‌ها را ساده می‌کند.
این پژوهش، دربارهٔ تحلیل شکل است، نه بازیابی جریان سیال.
بنابراین، معیار هندسی آن، زمان فیزیکی حرکت را تعیین نمی‌کند.
این تمایز، برای تفسیر مسیرهای بدون زمان نیز اهمیت دارد.
\cite{Srivastava2011ElasticCurves}

در مسئلهٔ کانال، مختصات مطلق مسیر نیز اطلاعات فیزیکی دارند.
فاصله از دیواره و مکان شروع، بر حرکت ذره اثر می‌گذارند.
پس حذف انتقال، دوران یا اندازهٔ منحنی، عموماً مجاز نیست.
چنین تبدیل‌هایی ممکن است مقایسهٔ ظاهری دو مسیر را بهتر کنند.
اما مسئلهٔ فیزیکی دیگری را جایگزین مسئلهٔ اولیه می‌کنند.
در این مقایسه، پارامتردهی حرکت می‌تواند مجهول باقی بماند.
مکان و طول مسیر، در دستگاه مختصات کانال حفظ می‌شوند.

\textbf{طول قوس و زمان‌بندی پنهان}

کمیت‌ها، مطابق رابطهٔ \ref{eq:lit-reference-scales}، بی‌بعد در نظر گرفته می‌شوند.
طول قوس نیز بر مقیاس طول $L_0$ تقسیم می‌شود.
منحنی ذرهٔ $i$ را با $\boldsymbol\gamma_i(s)$ نمایش می‌دهیم.
در اینجا، $s$ طول قوس بی‌بعد و $\ell_i$ طول منحنی است.
روی بخش‌های منظم، اندازهٔ مشتق منحنی نسبت به $s$ واحد است.
یک تابع افزایشی، طول قوس را به زمان پنهان متصل می‌کند.
\begin{equation}
\begin{gathered}
0\le s\le\ell_i,\qquad
\|\boldsymbol\gamma_i'(s)\|=1,\qquad
h_i(s)=\frac{\mathrm dt_i}{\mathrm ds}>0,\\
t_i(s)=t_{0,i}+\int_0^s h_i(\sigma)\,\mathrm d\sigma,\qquad
\mathbf v_{p,i}(s)=\frac{\boldsymbol\gamma_i'(s)}{h_i(s)},\\
\frac{\mathrm d\mathbf v_{p,i}}{\mathrm dt}
=\frac{\boldsymbol\gamma_i''(s)}{h_i(s)^2}
-\frac{\boldsymbol\gamma_i'(s)h_i'(s)}{h_i(s)^3}.
\end{gathered}
\label{eq:lit-untimed-clock}
\end{equation}
علامت پرایم، مشتق نسبت به طول قوس بی‌بعد را نشان می‌دهد.
مثبت‌بودن $h_i$، ترتیب مشاهده‌شدهٔ نقاط را در حرکت حفظ می‌کند.
همچنین، این رابطه نشان می‌دهد هندسه، سرعت و شتاب را تعیین نمی‌کند.
برای تعیین آن‌ها، زمان‌بندی نیز لازم است.
ازاین‌رو، تخمین جریان و تخمین زمان‌بندی به یکدیگر وابسته می‌شوند.

اگر سرعت اولیه معلوم باشد، بخشی از زمان‌بندی مقید می‌شود.
برای سرعت اولیهٔ غیرصفر، اندازه و جهت باید سازگار باشند.
برای نمونه، یک نمایش مثبت چنین نوشته می‌شود.
\begin{equation}
\begin{gathered}
h_i(0)=\frac{1}{\|\mathbf v_{p,0,i}\|},\qquad
\boldsymbol\gamma_i'(0)=
\frac{\mathbf v_{p,0,i}}{\|\mathbf v_{p,0,i}\|},\\
h_i(s)=h_i(0)\exp\!\left[g_i(s)-g_i(0)\right],\qquad
T_i=\int_0^{\ell_i}h_i(s)\,\mathrm ds.
\end{gathered}
\label{eq:lit-untimed-positive-clock}
\end{equation}
تابع $g_i$ می‌تواند با تعداد محدودی ضریب اسپلاین نمایش داده شود.
این نمایش، تنها نمونه‌ای برای اعمال مثبت‌بودن زمان‌بندی است.
مدت حرکت $T_i$ نیز از زمان‌بندی تخمینی حاصل می‌شود.
بدون مشاهدهٔ مستقل، نباید این مدت را معلوم فرض کرد.
اگر سرعت اولیه صفر باشد، رابطهٔ ابتدایی بالا کاربرد ندارد.
توقف ذره نیز به برخورد مناسب با تکینگی پارامتردهی نیاز دارد.
در چنین حالت‌هایی، صورت‌بندی باید برای بخش‌های منظم بازتعریف شود.

برای ذرهٔ اینرسی‌دار، مدل حرکت همچنان باید برقرار بماند.
رابطهٔ انتگرالی بخش قبل، برای فرض‌های درگ خطی قابل استفاده است.
در حالت بدون زمان، زمان‌های ورود به آن رابطه نیز مجهول‌اند.
پس منحنی مشاهده‌شده، همراه یک ساعت پنهان وارد برازش می‌شود.
زمان‌بندی دلخواه، لزوماً با دینامیک ذره و میدان مجاز سازگار نیست.
قیود فیزیکی و اطلاعات اولیه، مجموعهٔ زمان‌بندی‌های ممکن را محدود می‌کنند.
این محدودیت، به‌تنهایی اثبات یکتایی جریان یا زمان‌بندی نیست.

\textbf{حذف نشانه‌های زمانی از داده}

حذف ستون زمان، ممکن است اطلاعات زمانی را کاملاً حذف نکند.
اگر نقاط با فاصلهٔ زمانی ثابت ثبت شده باشند، تراکم مکانی معنادار است.
بخش‌های آهسته، معمولاً نقاط بیشتری در واحد طول خواهند داشت.
در نتیجه، فاصلهٔ نقاط می‌تواند نشانه‌ای از سرعت قبلی باقی بگذارد.
برای کاهش این نشانه، منحنی برحسب طول قوس بازنمونه‌برداری می‌شود.
درون‌یابی، هموارسازی و انتخاب ابتدا و انتها نیز باید ثبت شوند.
اگر مدت حرکت معلوم بماند، مسئله واقعاً فاقد همهٔ اطلاعات زمانی نیست.
پس اطلاعات حذف‌شده و اطلاعات باقی‌مانده باید جداگانه مشخص شوند.

در جریان پایا، مبدأ قراردادی زمان معمولاً مسئله را تغییر نمی‌دهد.
اما مدت پیمایش همچنان بر تفسیر دینامیکی اثر دارد.
در جریان ناپایا، زمان شروع می‌تواند فاز تحریک را تعیین کند.
بنابراین، انتقال مبدأ زمان همیشه یک انتخاب بی‌اثر نیست.
در مدل سینماتیکی ردیاب پایا، مقیاس سرعت و هندسه نیز متمایزند.
ضرب میدان سرعت در یک ثابت مثبت، خطوط جریان را تغییر نمی‌دهد.
این مشاهده، مربوط به سینماتیک ردیاب است.
آن را نباید ناوردایی عمومی معادلات جریان و ذره دانست.

\textbf{معیارهای هندسی و تفسیر مقایسه}

ارزیابی بدون زمان، به معیارهایی مستقل از ساعت اندازه‌گیری نیاز دارد.
یک معیار، اختلاف نقاط متناظر در طول قوس نرمال‌شده است.
معیار مکمل، فاصلهٔ هاسدورف میان مجموعهٔ نقاط دو منحنی است.
برای دو منحنی منظم، تعریف‌های زیر قابل استفاده‌اند.
\begin{equation}
\begin{aligned}
E_{\mathrm{arc}}^2={}&
\int_0^1\left\|
\boldsymbol\gamma_{\mathrm{obs}}(\ell_{\mathrm{obs}}\xi)
-\boldsymbol\gamma_{\mathrm{rep}}(\ell_{\mathrm{rep}}\xi)
\right\|^2\,\mathrm d\xi,\\
d_H(C_{\mathrm{obs}},C_{\mathrm{rep}})={}&
\max\left\{
\sup_{\mathbf a\in C_{\mathrm{obs}}}\inf_{\mathbf b\in C_{\mathrm{rep}}}
\|\mathbf a-\mathbf b\|,
\sup_{\mathbf b\in C_{\mathrm{rep}}}\inf_{\mathbf a\in C_{\mathrm{obs}}}
\|\mathbf b-\mathbf a\|
\right\}.
\end{aligned}
\label{eq:lit-untimed-geometric-errors}
\end{equation}
نماد $C$، مجموعهٔ نقاط منحنی مربوط را نشان می‌دهد.
پارامتر $\xi$، پیشروی نسبی روی طول منحنی است، نه زمان.
در معیار نخست، شروع و جهت پیمایش باید یکسان تعریف شوند.
فاصلهٔ هاسدورف، ترتیب عبور از نقاط را لحاظ نمی‌کند.
بنابراین، این دو معیار، اطلاعات کاملاً یکسانی ارائه نمی‌دهند.
هیچ‌کدام نیز خطای زمان حرکت را اندازه نمی‌گیرد.
زمان تخمینی، فقط با وجود مرجع زمانی مستقل قابل ارزیابی است.
پایان زودهنگام بازپخش یا بخش‌های مفقود مسیر، جداگانه گزارش می‌شوند.
نرمال‌کردن طول قوس، نباید ناقص‌بودن مسیر را پنهان کند.

برای مقایسهٔ دو روش، معیار هندسی باید مشترک باشد.
بهتر است داده، نویز و شرایط اولیه نیز یکسان باشند.
بودجهٔ بهینه‌سازی و اطلاعات جانبی، باید در مقایسه ثبت شوند.
اگر اجراها متفاوت باشند، اختلاف خطا فقط اثر حذف زمان نیست.
در پژوهش حاضر، جزئیات این صورت‌بندی در فصل سوم مشخص می‌شود.
نتایج بازپخش و بازیابی پروفایل نیز در فصل چهارم ارائه می‌شوند.

\زیرقسمت{استفاده از عملگر آموزش‌دیده در حل معکوس}
\label{sec:lit-frozen-operator-inverse}

\textbf{تفکیک آموزش مستقیم و تخمین ورودی}

حل معکوس، معمولاً به ارزیابی مکرر پاسخ مستقیم نیاز دارد.
اگر هر ارزیابی پرهزینه باشد، بهینه‌سازی نیز زمان‌بر خواهد شد.
یک عملگر آموزش‌دیده، می‌تواند این ارزیابی‌ها را تقریب بزند.
در \lr{DeepONet}، ورودی تابعی به پاسخ میدان نگاشت می‌شود.
این نگاشت، پس از آموزش برای ورودی‌های مجاز قابل ارزیابی است.
\cite{lu2021deeponet}

آموزش عملگر می‌تواند از داده، معادلات فیزیکی یا ترکیب آن‌ها استفاده کند.
در آموزش فیزیک‌آگاه، باقیماندهٔ معادلات برای خانواده‌ای از ورودی‌ها کنترل می‌شود.
اما آموزش چنین خانواده‌ای، با حل یک مسئلهٔ معکوس خاص متفاوت است.
مرحلهٔ نخست، وزن‌های شبکه را تعیین می‌کند.
مرحلهٔ دوم، ورودی سازگار با مشاهدهٔ جدید را جست‌وجو می‌کند.
\cite{wang2021pideeponet}

وزن‌های آموزش‌یافته را با $\theta^\star$ و مجهول‌های ورودی را با $\boldsymbol\alpha$ نشان می‌دهیم.
این دو مجموعه، نقش یکسانی در بازسازی ندارند.
برای مسئلهٔ ذره، پاسخ میدان باید به مدل مشاهده متصل شود.
یک نمایش کلی از این زنجیره، چنین است.
\begin{equation}
\begin{gathered}
\widehat{\mathbf q}_{\boldsymbol\alpha}
=\widehat{\mathcal G}_{\theta^\star}(\boldsymbol\alpha),
\qquad
\widehat{\mathbf q}_{\boldsymbol\alpha}=(\widehat{\mathbf u}_{\boldsymbol\alpha},
\widehat p_{\boldsymbol\alpha}),\\
\widehat{\mathbf d}(\boldsymbol\alpha,\boldsymbol\beta)
=\mathcal H_{\boldsymbol\beta}
\!\left[\mathcal P(\widehat{\mathbf u}_{\boldsymbol\alpha};\mathbf z_0)\right],
\qquad \theta=\theta^\star.
\end{gathered}
\label{eq:lit-frozen-observation-chain}
\end{equation}
نماد $\mathcal P$، نگاشت حرکت ذره با شرایط اولیهٔ $\mathbf z_0$ است.
نماد $\mathcal H$، نحوهٔ استخراج مشاهده از پاسخ را مشخص می‌کند.
بردار $\boldsymbol\beta$، در صورت نیاز، پارامترهای زمان‌بندی پنهان را دربرمی‌گیرد.
برای دادهٔ زمان‌دار، ساعت مشاهده معلوم است و این مجهول لازم نیست.
این زنجیره، مدل بازپخش مستقیم را نمایش می‌دهد.
برازش انتگرالی، می‌تواند مستقیم روی مسیر مشاهده‌شده تعریف شود.
در آن حالت، تعریف نگاشت مشاهده باید متناسب اصلاح شود.
این نمایش کلی، جایگزین صورت‌بندی دقیق تابع هدف نمی‌شود.

ثابت‌بودن وزن‌ها، به معنی ثابت‌بودن پاسخ نسبت به ورودی نیست.
در روش گرادیانی، مشتق پاسخ نسبت به $\boldsymbol\alpha$ همچنان لازم است.
بنابراین، گرادیان از شبکهٔ ثابت نیز عبور می‌کند.
در این مرحله، وزن‌های شبکه به‌روزرسانی نمی‌شوند.
تعداد مجهول‌های فیزیکی نیز با تعداد وزن‌های شبکه برابر نیست.
همچنین، پیشین ناشی از آموزش، دادهٔ اندازه‌گیری‌شدهٔ تازه محسوب نمی‌شود.

\textbf{روش‌های متفاوت استفاده از پیش‌آموزش}

تیواری و سن، عملگر عصبی را در طراحی معکوس بیزی به‌کار گرفتند.
مسئلهٔ آنان، تخمین هندسهٔ نازل از مشاهدات جریان تراکم‌پذیر است.
مدل مستقیم، بر معادلات اویلر شبه‌یک‌بعدی استوار است.
یک \lr{DeepONet}، نگاشت هندسه به توزیع عدد ماخ را تقریب می‌زند.
این جانشین، هنگام استنتاج ثابت می‌ماند و داخل درست‌نمایی ارزیابی می‌شود.
در کنار آن، یک شبکهٔ معکوس مستقیم نیز بررسی شده است.
این دو مسیر محاسباتی، با یکدیگر یکسان نیستند.
مطالعه، نمونه‌ای از استفادهٔ ثابت از عملگر مستقیم فراهم می‌کند.
اما هندسه، مشاهدات و فیزیک آن با مسئلهٔ کانال حاضر متفاوت‌اند.
\cite{Tiwari2026BayesianCFD}

روی سارکار و همکاران، پیش‌آموزش را به شکل دیگری استفاده کردند.
در چارچوب شناسایی پارامتر، ابتدا پاسخ از مشاهدات حسگرها آموخته می‌شود.
سپس، مدل پیش‌آموزش‌دیده همراه شبکهٔ تخمین پارامتر، آموزش را ادامه می‌دهد.
قید معادله و داده‌های موجود، این مرحله را هدایت می‌کنند.
بنابراین، این روش، صرفاً بهینه‌سازی ورودی یک عملگر ثابت نیست.
ازاین‌رو، عبارت پیش‌آموزش‌دیده به‌تنهایی برای تشخیص روش کافی نیست.
باید مشخص شود در مرحلهٔ معکوس، کدام پارامترها تغییر می‌کنند.
\cite{RoySarkar2026HiddenPhysics}

نمونهٔ دیگری از تطبیق، در چارچوب \lr{PINO} دیده می‌شود.
لی و همکاران، آموزش عملگر و پالایش پاسخ مسئلهٔ خاص را ترکیب کردند.
این پالایش، از قیود فیزیکی همان نمونه استفاده می‌کند.
پس هزینه و رفتار آن، با ارزیابی صرف یک جانشین ثابت متفاوت است.
این چارچوب مبتنی بر عملگر فوریه است و \lr{DeepONet} محسوب نمی‌شود.
بااین‌حال، برای تفکیک یادگیری خانواده و تطبیق نمونه، مقایسه‌ای مرتبط فراهم می‌کند.
\cite{Li2024PINO}

در روش معکوس مستقیم، نگاشت مشاهده به مجهول از پیش آموزش می‌بیند.
هنگام استفاده، این شبکه می‌تواند یک تخمین مستقیم تولید کند.
در وارون‌سازی عملگر مستقیم، برای هر مشاهده جست‌وجوی ورودی انجام می‌شود.
پس خروجی نهایی، حاصل ارزیابی یک شبکهٔ معکوس آموزش‌یافته نیست.
این تفاوت، بر هزینه، انعطاف و نحوهٔ اعتبارسنجی اثر می‌گذارد.
تخمین تک‌نقطه‌ای نیز لزوماً عدم‌قطعیت مجهول را توصیف نمی‌کند.

\textbf{نمایش مجهول و منظم‌سازی}

برای تخمین یک تابع، معمولاً نمایش محدودی از آن انتخاب می‌شود.
برای نمونه، پروفایل ورودی می‌تواند با ضرایب اسپلاین نمایش داده شود.
مقیاس سرعت نیز می‌تواند یک مجهول جداگانه باشد.
تعداد ضرایب، انعطاف خانواده و دشواری تخمین را تغییر می‌دهد.
افزایش ضرایب، لزوماً اطلاعات بیشتری به داده اضافه نمی‌کند.
همچنین، حدود مجاز باید با دامنهٔ اعتبار عملگر سازگار باشند.
حدود عددی هر مؤلفه، تمام ساختار خانوادهٔ آموزش را توصیف نمی‌کنند.
یک ترکیب تازه، ممکن است از الگوهای آموزشی فاصله داشته باشد.

برای نمایش ساختار کلی، می‌توان مسئله را چنین نوشت.
\begin{equation}
\begin{aligned}
(\widehat{\boldsymbol\alpha},\widehat{\boldsymbol\beta})
\in\underset{\boldsymbol\alpha\in\mathcal A,\,
\boldsymbol\beta\in\mathcal B}{\operatorname{argmin}}\quad
&\mathcal L_{\mathrm{obs}}
(\boldsymbol\alpha,\boldsymbol\beta;\theta^\star)\\
&+\lambda_s\mathcal R_s(\boldsymbol\alpha)
+\lambda_p\mathcal R_p(\boldsymbol\alpha)
+\lambda_t\mathcal R_t(\boldsymbol\beta)
+\lambda_f\mathcal R_f(\boldsymbol\alpha;\theta^\star).
\end{aligned}
\label{eq:lit-frozen-regularized-inverse}
\end{equation}
خطای $\mathcal L_{\mathrm{obs}}$، ناسازگاری با دادهٔ واقعی را اندازه می‌گیرد.
جملهٔ $\mathcal R_s$ می‌تواند نوسان یا خمیدگی پروفایل را جریمه کند.
جملهٔ $\mathcal R_p$، پیش‌دانستهٔ شکل یا مقیاس سرعت را بیان می‌کند.
جملهٔ $\mathcal R_t$، در حالت بدون زمان، زمان‌بندی را منظم می‌کند.
جملهٔ $\mathcal R_f$ می‌تواند باقیماندهٔ فیزیکی میدان را کنترل کند.
این کنترل، بهتر است به نقاط مسیر محدود نباشد.
ضرایب وزن، نقش داده و پیش‌دانسته را در تخمین تنظیم می‌کنند.
انتخاب جمله‌ها، به مسئلهٔ واقعی وابسته است.
این رابطه، الزام استفادهٔ هم‌زمان از تمام جریمه‌ها را ایجاد نمی‌کند.

همواری، یک پیش‌فرض دربارهٔ مجهول است، نه نتیجهٔ مستقیم مشاهده.
منظم‌سازی شدید می‌تواند جزئیات واقعی پروفایل را حذف کند.
منظم‌سازی ضعیف نیز ممکن است نوسان‌های سازگار با نویز ایجاد کند.
در مسئلهٔ غیرخطی، تخمین به مقداردهی اولیه حساس است.
چند شروع، امکان بررسی پاسخ‌های محلی متفاوت را فراهم می‌کند.
اما یافتن بهترین پاسخ میان شروع‌ها، اثبات بهینگی سراسری نیست.
پراکندگی این پاسخ‌ها نیز، به‌تنهایی توزیع پسین معتبر تولید نمی‌کند.

\textbf{خطای جانشین و اعتمادپذیری تخمین}

جانشین سریع، معمولاً پاسخ دقیق مدل مستقیم نیست.
کائو و همکاران، انتقال این خطا به استنتاج بیزی را بررسی کردند.
آنان نشان دادند خطای عملگر می‌تواند تخمین توزیع پسین را مختل کند.
برای اصلاح، یک مسئلهٔ تغییراتی خطی بر پایهٔ باقیمانده حل می‌شود.
کاهش مرتبهٔ خطا، به فرض‌های مشخص دربارهٔ مدل و تقریب وابسته است.
آزمون‌ها، شامل واکنش–پخش غیرخطی و مادهٔ فوق‌کشسان هستند.
این اصلاح، با افزودن جریمهٔ مومنتوم به تابع هدف یکسان نیست.
همچنین، این مطالعه صحت هر جانشین جریان را تضمین نمی‌کند.
\cite{Cao2023ResidualCorrection}

اثر خطای جانشین را می‌توان با یک تجزیهٔ ساده توضیح داد.
فرض می‌کنیم $\mathcal F$، مدل مستقیم مشاهده و $\boldsymbol\alpha_\dagger$ مجهول واقعی است.
برای دادهٔ زمان‌دار، مدل مشاهده در این رابطه ثابت است.
\begin{equation}
\begin{gathered}
\mathbf d=\mathcal F(\boldsymbol\alpha_\dagger)+\boldsymbol\varepsilon,\qquad
\mathbf e_\theta(\boldsymbol\alpha)
=\mathcal F(\boldsymbol\alpha)
-\widehat{\mathcal F}_{\theta^\star}(\boldsymbol\alpha),\\
\mathbf d-\widehat{\mathcal F}_{\theta^\star}(\boldsymbol\alpha_\dagger)
=\boldsymbol\varepsilon+\mathbf e_\theta(\boldsymbol\alpha_\dagger).
\end{gathered}
\label{eq:lit-frozen-surrogate-discrepancy}
\end{equation}
نماد $\boldsymbol\varepsilon$، خطای مشاهده نسبت به مدل مرجع را نشان می‌دهد.
نماد $\mathbf e_\theta$، اختلاف جانشین و همان مدل مرجع است.
بنابراین، حتی ورودی واقعی نیز ممکن است خطای برازش غیرصفر داشته باشد.
بهینه‌ساز می‌تواند برای جبران خطای جانشین، ورودی را تغییر دهد.
در این صورت، برازش بهتر الزاماً بازیابی بهتر مجهول نیست.
همچنین، دقت خروجی، دقت مشتق نسبت به ورودی را تضمین نمی‌کند.
این موضوع، برای جست‌وجوی گرادیانی اهمیت دارد.
ارزیابی جانشین باید ناحیهٔ ورودی‌های محتملِ تخمین را نیز پوشش دهد.

\textbf{بازپخش، هزینه و مقایسهٔ روش‌ها}

پس از برازش، میدان متناظر با مجهول تخمینی محاسبه می‌شود.
سپس، حرکت ذره از شرایط اولیه دوباره حل می‌شود.
در بازپخش، میدان در مکان در حال تغییر ذره ارزیابی می‌شود.
برای مسیر بدون زمان، ابتدا مقایسهٔ هندسی مشترک انجام می‌شود.
اگر زمان مرجع موجود باشد، ساعت تخمینی نیز جداگانه سنجیده می‌شود.
بازپخش با حل‌گر مستقل، خطای انتگرال‌گیری را کنترل می‌کند.
اما مرجع مستقل میدان، آزمونی متفاوت برای خطای جانشین است.
در مسائل مصنوعی، مقایسهٔ پروفایل واقعی و تخمینی نیز ضروری است.
بدون مرجع مجهول، تنها سازگاری با داده و قیود قابل گزارش است.

هزینهٔ تولید داده و آموزش، پیش از بازسازی پرداخت می‌شود.
هزینهٔ هر تخمین، شامل ارزیابی‌ها، مشتق‌ها و تکرارهای جست‌وجو است.
بازپخش و کنترل با حل‌گر مرجع نیز هزینهٔ جداگانه دارند.
مقایسهٔ زمانی، باید دقت هدف و سخت‌افزار را مشخص کند.
برای یک بازسازی، هزینهٔ آموزش ممکن است غالب باشد.
برای بازسازی‌های تکراری، این هزینه میان مسائل بیشتری توزیع می‌شود.
پس سرعت ارزیابی شبکه، به‌تنهایی شتاب کل فرایند را تعیین نمی‌کند.

جدول \ref{tab:lit-inverse-operator-strategies} تفاوت نقش شبکه در چهار راهبرد را نشان می‌دهد.
این راهبردها، پاسخ یکسانی به پرسشِ مجهول‌بودن وزن‌ها نمی‌دهند.

\begin{table}[htbp]
\centering
\small
\caption{تفاوت راهبردهای استفاده از شبکهٔ آموزش‌دیده در حل معکوس}
\label{tab:lit-inverse-operator-strategies}
\begin{tabular}{p{0.24\textwidth}p{0.30\textwidth}p{0.35\textwidth}}
\hline
راهبرد & محاسبه در مرحلهٔ استفاده & نکتهٔ تفکیک \\
\hline
وارون‌سازی عملگر مستقیم ثابت
& جست‌وجوی ورودی‌های فیزیکی
& وزن‌ها ثابت؛ نمونهٔ تیواری و سن \cite{Tiwari2026BayesianCFD} \\
تطبیق فیزیک‌آگاه مدل
& به‌روزرسانی مدل و مجهول‌ها
& ادامهٔ آموزش؛ نمونهٔ روی سارکار \cite{RoySarkar2026HiddenPhysics} \\
شبکهٔ معکوس مستقیم
& ارزیابی نگاشت مشاهده به مجهول
& نگاشت معکوس از پیش آموخته شده است \\
اصلاح پاسخ جانشین
& حل مسئلهٔ کمکی برای پاسخ میدان
& اصلاح تغییراتی؛ نمونهٔ کائو \cite{Cao2023ResidualCorrection} \\
\hline
\end{tabular}
\end{table}

در پژوهش حاضر، عملگر مستقیم در مرحلهٔ تخمین ثابت نگه داشته می‌شود.
مجهول‌های ورودی، و در حالت لازم زمان‌بندی، با مشاهدات سازگار می‌شوند.
این انتخاب، نقش مدل مستقیم و دادهٔ محدود را روشن می‌کند.
اما ادعای بازیابی، باید با فرض‌ها و دامنهٔ اعتبار همراه باشد.
صورت‌بندی اجرایی در فصل سوم و ارزیابی مستقل در فصل چهارم می‌آیند.
در ادامه، تخمین تحریک مغناطیسی و محدودیت‌های شناسایی‌پذیری بررسی می‌شوند.

\زیرقسمت{تخمین تحریک مغناطیسی و رهگیری مسیر}
\label{sec:lit-magnetic-inverse-tracking}

\textbf{دو پرسش متفاوت دربارهٔ تحریک}

در هدایت مغناطیسی، می‌توان فرمانی برای رسیدن به مسیر هدف انتخاب کرد.
در این حالت، مسیر مطلوب داده شده و فرمان طراحی می‌شود.
پرسش دیگر، بازیابی فرمانی است که حرکت مشاهده‌شده را ایجاد کرده است.
در این مسئله، دادهٔ اصلی حرکت ذره و مجهول، تحریک گذشته است.
این دو پرسش، الزاماً پاسخ یا معیار موفقیت یکسانی ندارند.
یک فرمان مناسب برای رهگیری، ممکن است فرمان تاریخی واقعی نباشد.
پس نزدیکی مسیرها، برای اثبات صحت بازیابی تحریک کافی نیست.
پژوهش حاضر، طراحی فرمان برای بازتولید مسیر هدف را دنبال می‌کند.

الحمداوی و همکاران، رهگیری ذرهٔ مغناطیسی را با قیود حالت بررسی کردند.
کنترل در سطح نیروی مغناطیسی تعریف شده است.
کنترلگرهای خطی و همگن، در حضور اغتشاش محدود مقایسه می‌شوند.
ارزیابی گزارش‌شده، مبتنی بر شبیه‌سازی است.
این مطالعه، وارون‌سازی آرایش واقعی آهنرباها برای تولید هر نیرو را حل نمی‌کند.
بنابراین، رهگیری در سطح نیرو، با بازیابی فرمان سخت‌افزاری یکسان نیست.
این تمایز، در تفسیر الگوریتم‌های هدایت ذره اهمیت دارد.
\cite{ElHamdaoui2026Tracking}

استرایر و همکاران، توزیع نیروی مغناطیسی را از حرکت ذرات برآورد کردند.
آنان، ریزذرات پلی‌استایرن را در سیالی با پاسخ پارامغناطیسی دنبال کردند.
سرعت‌های حاصل از رهگیری، همراه خواص معلوم، وارد توازن نیرو می‌شوند.
اثر وزن و شناوری نیز در محاسبه لحاظ می‌شود.
سپس، توزیع مکانی تخمینی با محاسبات سه‌بعدی مقایسه می‌شود.
این کار، نمونه‌ای تجربی از استنتاج نیروی محرک از حرکت است.
اما مجهول آن، دو فرمان زمانی مستقل برای دو منبع نیست.
\cite{Strayer2024ForceMapping}

\textbf{مدل میدان و دادهٔ قابل استنتاج}

مدل انتقال مغناطیسی، باید حرکت سیال و پاسخ ذره را مشخص کند.
مطالعهٔ فورلانی، نمونه‌ای از مدل‌سازی انتقال در یک ریزسامانهٔ مغناطیسی است.
این چارچوب‌های فیزیکی، مبنای تعریف رابطهٔ ورودی و حرکت قرار می‌گیرند.
بااین‌حال، وجود مدل مستقیم، یکتایی وارون آن را تضمین نمی‌کند.
\cite{Furlani2006Transport}

برای مقایسه، مدل سادهٔ دو منبع ثابت را در نظر می‌گیریم.
مقیاس‌های طول، سرعت و زمان، مطابق بخش‌های قبل ثابت هستند.
میدان نیز بر مقدار مرجع مثبت $B_0$ تقسیم می‌شود.
پس مکان، زمان، سرعت و میدان در روابط بعدی بی‌بعد هستند.
برای پاسخ القایی خطی و حرکت بیش‌میرا، می‌نویسیم:
\begin{equation}
\begin{gathered}
\mathbf b(\mathbf r,t)
=a_L(t)\mathbf b_L(\mathbf r)
+a_R(t)\mathbf b_R(\mathbf r),\\
\frac{\mathrm d\mathbf r}{\mathrm dt}
=\mathbf u(\mathbf r,t)
+\Lambda_m\nabla_{\mathbf r}|\mathbf b(\mathbf r,t)|^2,
\qquad
\Lambda_m=\frac{\alpha_m B_0^2}{2\zeta U_0L_0}.
\end{gathered}
\label{eq:lit-magnetic-dimensionless-observation}
\end{equation}
توابع $\mathbf b_L$ و $\mathbf b_R$، شکل فضایی معلوم دو منبع هستند.
دامنه‌های $a_L$ و $a_R$، فرمان‌های بی‌بعد منابع را نشان می‌دهند.
ضریب $\alpha_m$، پاسخ مغناطیسی خطی ذره را مشخص می‌کند.
ضریب $\zeta$، مقاومت انتقالی بُعددار ذره است.
در این مدل، $\Lambda_m$ معلوم و ثابت فرض می‌شود.
همچنین، اینرسی ذره و حافظهٔ مغناطیسی نادیده گرفته می‌شوند.
حرکت براونی و تصحیح درگ دیواره، در این تقریب لحاظ نمی‌شوند.
کاربرد این رابطه، به اعتبار همین فرض‌ها وابسته است.

حرکت، ابتدا سهم سرعت ناشی از نیروی مغناطیسی را آشکار می‌کند.
تبدیل آن به فرمان منابع، یک مرحلهٔ وارون دیگر است.
یک تحریک مؤثر، الزاماً دو دامنهٔ مستقل را تعیین نمی‌کند.
اگر جریان حامل مجهول باشد، تفکیک دو سهم دشوارتر می‌شود.
عدم‌قطعیت خواص مغناطیسی نیز می‌تواند با دامنهٔ تحریک اشتباه شود.
پس معلوم‌بودن جریان و خواص ذره، اطلاعات جانبی مهمی فراهم می‌کند.

بسط عبارت میدان، ساختار این وارون‌سازی را روشن می‌کند.
\begin{equation}
\begin{aligned}
\nabla|\mathbf b|^2={}&
a_L^2\nabla|\mathbf b_L|^2
+a_R^2\nabla|\mathbf b_R|^2\\
&+2a_La_R\nabla(\mathbf b_L\cdot\mathbf b_R),\\
\nabla|\mathbf b(-a_L,-a_R)|^2={}&
\nabla|\mathbf b(a_L,a_R)|^2.
\end{aligned}
\label{eq:lit-magnetic-sign-symmetry}
\end{equation}
نیرو، شامل دو سهم مربعی و یک سهم متقاطع است.
با تغییر هم‌زمان علامت دو فرمان، میدان علامت عوض می‌کند.
اما عبارت مؤثر در حرکت، تغییری نمی‌کند.
برای شرایط اولیهٔ یکسان، مدل می‌تواند مسیر یکسانی تولید کند.
اگر هر دو فرمان در بازه‌ای صفر باشند، پالس‌ها جدا می‌شوند.
در مدل بدون حافظه، علامت پالس‌های جدا می‌تواند مستقل تغییر کند.
این تغییر، به مجازبودن فرمان‌ها و حفظ همواری لازم وابسته است.
بنابراین، ابهام ممکن است از یک علامت کلی گسترده‌تر باشد.
اگر هر دو فرمان مجاز باشند، این ابهام ساختاری باقی می‌ماند.
بنابراین، حرکت به‌تنهایی همیشه علامت مشترک فرمان‌ها را تعیین نمی‌کند.

اطلاع مستقل از قطبیت، ممکن است این دوگانگی را حذف کند.
محدودیت فیزیکی دامنه نیز باید در صورت‌بندی ثبت شود.
بااین‌حال، انتخاب قرارداد علامت، اثبات یکتایی کامل دو تابع نیست.
همچنین، تقارن بالا به همین مدل و نبود میدان زمینهٔ ثابت وابسته است.
نباید آن را به همهٔ سامانه‌های مغناطیسی تعمیم داد.

\textbf{نمایش فرمان و تابع هدف}

در تخمین تاریخی تحریک، مسیر زمان‌دار مشاهده می‌شود.
در طراحی، همان صورت تابع هدف برای بازتولید مسیر به کار می‌رود.
فرمان‌ها می‌توانند با تعداد محدودی تابع پایه نمایش داده شوند.
انتخاب تعداد ضرایب، قدرت تفکیک زمانی تخمین را محدود می‌کند.
این محدودیت باید با نرخ ثبت و دامنهٔ آموزش سازگار باشد.
برای یک مدل مستقیم ثابت، صورت کلی زیر قابل استفاده است.
\begin{equation}
\begin{gathered}
a_k(t)=\sum_{j=1}^{m_k}c_{kj}\phi_{kj}(t),
\qquad k\in\{L,R\},\\
\widehat{\mathbf c}\in
\underset{\mathbf c\in\mathcal C}{\operatorname{argmin}}
\left\{
\sum_n\|\widehat{\mathbf r}(t_n;\mathbf c)
-\mathbf r_n^{\mathrm{obs}}\|_{W_n}^2
+\lambda\sum_{k\in\{L,R\}}\int_0^T|\dot a_k(t)|^2\,\mathrm dt
\right\},\\
a_{k,\min}\le a_k(t)\le a_{k,\max}.
\end{gathered}
\label{eq:lit-magnetic-input-estimation}
\end{equation}
توابع $\phi_{kj}$، پایه‌های زمانی و ضرایب $c_{kj}$، مجهول‌های تخمین هستند.
مجموعهٔ $\mathcal C$، قیود دامنه و سایر محدودیت‌های مجاز را اعمال می‌کند.
ماتریس‌های $W_n$، وزن اختلاف موقعیت‌ها را مشخص می‌کنند.
جریمهٔ زمانی، فرمان‌های نوسانی را محدود می‌کند.
این جریمه، اطلاعات اندازه‌گیری‌شدهٔ تازه‌ای ایجاد نمی‌کند.
اگر سرعت مستقیماً اندازه‌گیری شود، می‌توان جملهٔ متناظر آن را افزود.
اما سرعت مشتق‌شده از موقعیت، مشاهده‌ای مستقل از همان موقعیت نیست.

در صورت استفاده از عملگر پنجره‌ای، اتصال پنجره‌ها نیز اهمیت دارد.
حالت انتهای پنجره، بر حرکت پنجرهٔ بعد اثر می‌گذارد.
برازش جداگانه، ممکن است سازگاری مسیر کامل را پنهان کند.
پس فرمان بازیابی‌شده باید در بازپخش کامل بررسی شود.
خطای مسیر و خطای فرمان نیز جداگانه گزارش می‌شوند.
جزئیات مدل و قیود واقعی، در فصل سوم مشخص خواهند شد.

\زیرقسمت{شناسایی‌پذیری، نویز و اعتبارسنجی}
\label{sec:lit-identifiability-noise-validation}

\textbf{برازش موفق و بازیابی یکتا}

شناسایی‌پذیری، به امکان تشخیص مجهول از مشاهدات مربوط است.
برازش موفق، فقط سازگاری یک پاسخ با داده را نشان می‌دهد.
ممکن است پاسخ‌های دیگری نیز مشاهدات مشابهی ایجاد کنند.
در شناسایی‌پذیری ساختاری، مدل و نوع مشاهده در شرایط ایدئال بررسی می‌شوند.
در شناسایی‌پذیری عملی، کیفیت و محدودیت داده نیز اهمیت دارد.
یک مجهول، ممکن است از نظر ساختاری قابل تشخیص باشد.
اما دادهٔ واقعی ممکن است برای تخمین دقیق کافی نباشد.
همچنین، یکتایی محلی با یکتایی در تمام دامنهٔ مجاز متفاوت است.

راوه و همکاران، از درست‌نمایی نمایه‌ای برای تحلیل پارامترها استفاده کردند.
در این روش، یک پارامتر ثابت و سایر پارامترها دوباره بهینه می‌شوند.
این جست‌وجو، جهت‌های ضعیفاً مقید و وابستگی پارامترها را آشکار می‌کند.
کاربرد اصلی مقاله، مدل‌های دینامیکی با مشاهدهٔ ناقص است.
بنابراین، این روش یک ابزار تشخیصی مرتبط فراهم می‌کند.
اما نتیجهٔ یک مدل، به مدل ذره و مشاهدهٔ دیگری منتقل نمی‌شود.
\cite{Raue2009ProfileLikelihood}

\textbf{نقش محل مشاهده و اطلاعات جانبی}

تعداد داده‌ها، تمام محتوای اطلاعاتی آن‌ها را مشخص نمی‌کند.
دو مسیر نزدیک، ممکن است نسبت به مجهول‌ها حساسیت مشابهی داشته باشند.
در مقابل، مسیرهای متفاوت می‌توانند ترکیب‌های تازه‌ای از پارامترها را آشکار کنند.
این موضوع، به مکان رهاسازی و شرایط اولیه نیز وابسته است.
افزایش مدت ثبت، لزوماً جهت‌های حساسیت تازه‌ای ایجاد نمی‌کند.
همبستگی زمانی داده‌ها نیز باید در تفسیر تعداد نمونه‌ها لحاظ شود.

وانگ و زکی، نقش محل مشاهده را در جریان آشفته بررسی کردند.
آنان، داده‌های سرعت ناحیهٔ بیرونی را برای بازسازی نزدیک دیواره استفاده کردند.
روش تغییراتی، تحول پاسخ را به معادلات ناویر–استوکس مقید می‌کند.
با دورشدن مشاهدات، بازیابی برخی مقیاس‌های نزدیک دیواره دشوارتر می‌شود.
اثر تفکیک مکانی و زمانی داده نیز بررسی شده است.
این مسئله، مشاهدهٔ اویلری دارد و با دو مسیر ذره یکسان نیست.
بااین‌حال، اهمیت محل داده را فراتر از شمار نمونه‌ها نشان می‌دهد.
\cite{WangZaki2025OuterObservations}

ژو و گراور، نقش نویز، تراکم و اینرسی را بررسی کردند.
در بازسازی آنان، برخی ویژگی‌های ذره نیز مجهول هستند.
\cite{ZhouGrauer2026Joint}

برای تخمین تحریک مغناطیسی، شکل فضایی منابع نیز اهمیت دارد.
اگر دو منبع اثرهای مشابهی روی مسیر داشته باشند، تفکیکشان دشوار می‌شود.
همچنین، بخش‌های کم‌حساسیت مسیر، اطلاعات اندکی دربارهٔ تغییر فرمان منتقل می‌کنند.
اگر ضریب پاسخ مغناطیسی مجهول باشد، ابهام مقیاس نیز ممکن است ایجاد شود.
در مدل مربعی، افزایش دامنه می‌تواند با کاهش ضریب پاسخ جبران شود.
بنابراین، کالیبراسیون خواص ذره و جریان حامل، بخشی از اطلاعات مسئله است.

پیش‌دانسته‌ها، مانند همواری و محدودیت دامنه، مجموعهٔ پاسخ‌ها را کوچک می‌کنند.
اما کوچک‌شدن این مجموعه، لزوماً حاصل اطلاعات جدید در مشاهده نیست.
گزارش تخمین باید سهم داده و سهم پیش‌دانسته را روشن کند.
این نکته برای عملگر پیش‌آموزش‌دیده نیز برقرار است.

\textbf{حساسیت محلی و درست‌نمایی نمایه‌ای}

برای تحلیل محلی، تمام مجهول‌های تخمین در بردار $\boldsymbol\beta$ جمع می‌شوند.
این بردار می‌تواند ضرایب فرمان، پروفایل یا زمان‌بندی را شامل شود.
نگاشت پیش‌بینی مشاهده را با $\mathcal F(\boldsymbol\beta)$ نمایش می‌دهیم.
مقیاس پارامترها و مدل خطای مشاهده، باید از ابتدا مشخص باشند.
برای نویز گاوسی با کوواریانس معلوم و ثابت، می‌توان نوشت:
\begin{equation}
J=\frac{\partial\mathcal F}{\partial\boldsymbol\beta},
\qquad
\widetilde J=\Sigma^{-1/2}J,
\qquad
I=J^{\mathsf T}\Sigma^{-1}J
=\widetilde J^{\mathsf T}\widetilde J.
\label{eq:lit-identifiability-sensitivity}
\end{equation}
ماتریس $\Sigma$، کوواریانس مثبت‌معین خطای مشاهده است.
ماتریس $I$، اطلاعات فیشر محلی برای همین فرض‌های آماری است.
مقادیر منفرد کوچک، جهت‌های کم‌حساسیت را در تقریب خطی نشان می‌دهند.
بااین‌حال، مقدار آن‌ها به مقیاس انتخاب‌شدهٔ پارامترها نیز وابسته است.
کمبود رتبه، ناکافی‌بودن اطلاعات مرتبهٔ اول را نشان می‌دهد.
این وضعیت، به‌تنهایی شناسایی‌ناپذیری مدل غیرخطی را اثبات نمی‌کند.
از سوی دیگر، رتبهٔ کامل، وجود پاسخ‌های دوردست را رد نمی‌کند.
تقارن علامت مشترک، نمونه‌ای از ابهام خارج تحلیل صرفاً محلی است.

برای بررسی فراتر از همسایگی جواب، می‌توان نمایهٔ درست‌نمایی ساخت.
اگر $\ell$ منفی لگاریتم درست‌نمایی باشد، تعریف کلی چنین است:
\begin{equation}
\ell_{\mathrm{prof},j}(c)
=\min_{\boldsymbol\beta:\,\beta_j=c}
\ell(\boldsymbol\beta).
\label{eq:lit-identifiability-profile}
\end{equation}
در هر مقدار $c$، سایر مجهول‌ها دوباره بهینه می‌شوند.
ناحیهٔ کم‌تغییر نمایه، می‌تواند ضعف تعیین آن پارامتر را نشان دهد.
کیفیت این نتیجه، به حل بهینه‌سازی و حدود جست‌وجو وابسته است.
نمایهٔ عددی تخت، به‌تنهایی اثبات دقیق ابهام ساختاری نیست.
همچنین، افزودن جریمهٔ دلخواه، تفسیر آماری درست‌نمایی را تغییر می‌دهد.
بازهٔ اطمینان، به مدل نویز و شرایط استنباط معتبر نیاز دارد.

\textbf{نویز، تفکیک زمانی و خطای مدل}

نویز موقعیت، هنگام مشتق‌گیری زمانی تقویت می‌شود.
سرعت و شتاب حاصل، به روش هموارسازی نیز وابسته‌اند.
صورت‌بندی انتگرالی می‌تواند نیاز به مشتق دوم عددی را کاهش دهد.
اما خطای موقعیت، شرایط اولیه و مدل فیزیکی همچنان باقی می‌مانند.
اگر سرعت از همان موقعیت ساخته شود، خطاهایشان همبسته هستند.
پس افزودن دو جملهٔ مستقل، می‌تواند وزن اطلاعات تکراری را افزایش دهد.

نرخ ثبت، باید با تغییرات قابل بازیابی تحریک سازگار باشد.
بدون پیش‌دانسته، تغییرات سریع ممکن است از داده قابل تفکیک نباشند.
افزایش تعداد گره‌های فرمان، این کمبود اطلاعات را برطرف نمی‌کند.
در عملگر پنجره‌ای، نمونه‌برداری ورودی نیز محدودیت دیگری ایجاد می‌کند.
فرمان‌هایی با نمونه‌های یکسان، ممکن است جزئیات میان‌حسگری متفاوتی داشته باشند.
پس تعمیم باید در خانوادهٔ زمانی قابل نمایش ارزیابی شود.

اثر خطای جانشین، در بخش \ref{sec:lit-frozen-operator-inverse} بررسی شد.
این خطا، باید از نویز اندازه‌گیری جدا بماند.
\cite{Cao2023ResidualCorrection}

اگر نیروی دیگری حذف شده باشد، فرمان ممکن است اثرش را جبران کند.
در این حالت، فرمان تخمینی الزاماً فرمان واقعی دستگاه نیست.
اعتبار مدل بیش‌میرا و مدل اینرسی‌دار نیز باید جدا بررسی شود.
تغییر مدل در بازپخش، باید همراه توضیح تفاوت فرض‌ها باشد.

\textbf{اعتبارسنجی در چند سطح}

نخست، سازگاری پاسخ با مشاهدات برازش‌شده سنجیده می‌شود.
سپس، حرکت ذره با شرایط اولیه و ورودی تخمینی بازپخش می‌شود.
بهتر است انتگرال‌گیری با پیاده‌سازی مستقل نیز کنترل شود.
اگر همان جانشین استفاده شود، این آزمون استقلال میدان را تأمین نمی‌کند.
مرجع عددی معتبر، مقایسهٔ دیگری برای ارزیابی میدان و حرکت فراهم می‌کند.
در دادهٔ مصنوعی، مجهول واقعی نیز باید جداگانه مقایسه شود.
در دادهٔ تجربی، کالیبراسیون و اندازه‌گیری مستقل ارزش ویژه‌ای دارند.

خطای بازپخش مسیر را می‌توان در مقیاس بی‌بعد تعریف کرد.
برای فرمان، نحوهٔ برخورد با تقارن علامت نیز باید مشخص باشد.
اگر هر دو علامت مجاز باشند، یک معیار تکمیلی چنین است:
\begin{equation}
\begin{gathered}
E_{\mathrm{rep}}^2
=\frac{1}{T}\int_0^T
\|\mathbf r_{\mathrm{rep}}(t)-\mathbf r_{\mathrm{ref}}(t)\|^2\,\mathrm dt,\\
E_{a,\pm}
=\min_{\sigma\in\{-1,1\}}
\left[
\frac{\displaystyle\sum_{k\in\{L,R\}}\int_0^T
|\widehat a_k(t)-\sigma a_{k,\mathrm{ref}}(t)|^2\,\mathrm dt}
{\displaystyle\sum_{k\in\{L,R\}}\int_0^T
|a_{k,\mathrm{ref}}(t)|^2\,\mathrm dt}
\right]^{1/2}.
\end{gathered}
\label{eq:lit-magnetic-validation-errors}
\end{equation}
در این تعریف، $T$ مدت مشاهدهٔ بی‌بعد و مثبت است.
علامت $\sigma$، برای هر دو منبع و تمام زمان‌ها مشترک است.
این معیار، فقط یک تقارن علامت سراسری را لحاظ می‌کند.
هم‌ارزی پالس‌های جدا، به بررسی زمانی مستقل نیاز دارد.
مخرج معیار فرمان، باید غیرصفر باشد.
برای فرمان مرجع نزدیک صفر، خطای مطلق نیز لازم است.
اگر قطبیت معلوم باشد، خطای فرمان بدون اصلاح علامت گزارش می‌شود.
خطای میدان یا نیرو نیز با خطای فرمان یکسان نیست.
در نبود مرجع واقعی، ادعای صحت بازیابی باید محدود بماند.

آزمون هم‌مدل، داده را با همان جانشین تولید و بازسازی می‌کند.
این آزمون، سازگاری داخلی روش را بررسی می‌کند.
اما خطاهای مشترک تولید و تخمین، ممکن است در آن پنهان بمانند.
آزمون مرجع مستقل، باید تفاوت محاسباتی خود را مشخص کند.
مدل، شبکه، گام زمانی و دقت حل، از این مشخصات هستند.
افزودن نویز معلوم و تکرار آزمون‌ها، پایداری عملی را بهتر آشکار می‌کند.
این تکرارها، جای اثبات یکتایی را نمی‌گیرند.

جدول \ref{tab:lit-inverse-information-comparison} تفاوت مشاهده، مجهول و مرجع پژوهش‌ها را نشان می‌دهد.
این مطالعات، یک آزمون مشترک برای رتبه‌بندی دقت تشکیل نمی‌دهند.
تعداد مسیرها نیز بدون توجه به این تفاوت‌ها قابل مقایسه نیست.

\begin{table}[p]
\centering
\small
\caption{مقایسهٔ مدل، داده، مجهول و ارزیابی در پژوهش‌های مرتبط}
\label{tab:lit-inverse-information-comparison}
\setlength{\tabcolsep}{3pt}
\begin{tabular}{p{0.13\textwidth}p{0.15\textwidth}p{0.15\textwidth}p{0.15\textwidth}p{0.14\textwidth}p{0.15\textwidth}}
\hline
پژوهش & مدل فیزیکی & دادهٔ مشاهده & مجهول یا کنترل & پیش‌دانستهٔ مهم & مرجع ارزیابی \\
\hline
وان و همکاران \cite{Wan2026TrajectoryFlowNet}
& جریان و ردیاب غیرفعال
& موقعیت و سرعت زمان‌دار
& مسیر و میدان جریان
& فرض ردیاب و معادلات جریان
& شبیه‌سازی و دادهٔ تجربی \\
ژو و گراور \cite{ZhouGrauer2026Joint}
& جریان و دینامیک ذره
& موقعیت‌های زمان‌دار
& جریان و ویژگی‌های ذره
& قوانین حرکت و مدل عدم‌قطعیت
& مرجع عددی داده‌های مصنوعی \\
استرایر و همکاران \cite{Strayer2024ForceMapping}
& توازن نیرو در مغناطوفورز
& سرعت حاصل از رهگیری
& توزیع محرک مغناطیسی
& خواص ذره، سیال و شناوری
& مقایسه با حل سه‌بعدی \\
الحمداوی و همکاران \cite{ElHamdaoui2026Tracking}
& حرکت ذره با ورودی نیرو
& حالت و مسیر مرجع
& نیروی کنترل رهگیری
& مدل دینامیکی و حدود اغتشاش
& شبیه‌سازی کنترلگرها \\
وانگ و زکی \cite{WangZaki2025OuterObservations}
& ناویر–استوکس آشفته
& سرعت ناحیهٔ بیرونی
& حالت جریان و کمیت‌های دیواره
& معادلات و شرایط مرزی
& دادهٔ مرجع حل مستقیم \\
\hline
\end{tabular}
\end{table}

برای پژوهش حاضر، نوع مشاهده و مجهول باید در هر آزمایش ثبت شود.
تعداد مسیرها، اطلاعات زمانی و سرعت اولیه نیز صریح گزارش می‌شوند.
ادعای بازسازی، به خانوادهٔ پاسخ و فرض‌های مدل محدود می‌ماند.
در مسئلهٔ مغناطیسی، مسیر، نیرو و فرمان جداگانه ارزیابی می‌شوند.
معیار اصلی پژوهش حاضر، بازپخش مستقل مسیر هدف است.
مقایسهٔ نیرو در امتداد مسیر، یک بررسی تشخیصی است.
شیوهٔ اجرای این ارزیابی‌ها، در فصل‌های سوم و چهارم ارائه می‌شود.

\قسمت{تلفیق پیشینه و جایگاه پژوهش حاضر}
\label{sec:lit-review-synthesis}

پیشینهٔ بررسی‌شده، چهار مسیر مرتبط برای مطالعهٔ حرکت ذرات نشان می‌دهد.
طراحی رهاسازی، رفتار جمعی را از طریق شرایط شروع تنظیم می‌کند.
یادگیری عملگر، ارزیابی پاسخ ورودی‌های متفاوت را هدف می‌گیرد.
بازسازی جریان، حرکت ذره را به اطلاعات میدان پیوند می‌دهد.
تخمین تحریک، ورودی سازگار با حرکت مشاهده‌شده را جست‌وجو می‌کند.
در هر مسیر، تعریف داده و مجهول، مبنای مقایسه است.

\textbf{طراحی رهاسازی و حدود تشکیل آرایش}

در مطالعات هدایت، جریان یا تحریک می‌تواند هنگام حرکت تغییر کند.
تلهٔ استوکس و دست‌کاری اپتوسیالی، نمونه‌هایی از این مسیر هستند.
آنان، کنترل موقعیت و آرایش را با سازوکارهای فعال بررسی می‌کنند.
بنابراین، موفقیتشان مستقیماً عملکرد یک برنامهٔ رهاسازی ثابت را تعیین نمی‌کند.
\cite{Shenoy2016StokesTrap,Erben2024Optofluidic}

در صورت‌بندی حاضر، شرایط رهاسازی متغیرهای اصلی طراحی هستند.
شبکهٔ فیزیک‌آگاه پارامتریک، مسیر ذره را از شرایط رهاسازی پیش‌بینی می‌کند.
برنامهٔ رهاسازی با استفاده از این مسیرهای پیش‌بینی‌شده ساخته می‌شود.
مکان، سرعت و زمان رهاسازی، مجموعهٔ مسیرهای قابل‌دستیابی را تعیین می‌کنند.
تشکیل آرایش، با موقعیت مراکز در یک لحظهٔ مشترک سنجیده می‌شود.
حدود پرتابگرها و فاصلهٔ ذرات، بخشی از همین تعریف هستند.
حفظ آرایش پس از آن لحظه، پرسشی متفاوت است.
پس پیچیدگی شکل، باید همراه تعداد ذرات و منابع طراحی گزارش شود.
بازپخش مستقل نیز برای کنترل برنامهٔ منتخب اهمیت دارد.

پرسش این محور، دامنهٔ آرایش‌های قابل‌دستیابی با منابع محدود است.
افزایش ایستگاه‌ها، آزادی بیشتری ایجاد می‌کند، اما همگرایی عددی را تضمین نمی‌کند.
همچنین، قید فاصلهٔ مراکز جای مدل کامل برخورد را نمی‌گیرد.
ازاین‌رو، ادعای تشکیل آرایش به فرض‌های مدل و آزمون فاصله وابسته می‌ماند.

\textbf{یادگیری پاسخ برای ورودی‌های تازه}

در یادگیری عملگر، نوع تغییر ورودی باید مشخص باشد.
تغییر شکل پروفایل، تغییر دامنه و تغییر هندسه، مسائل متفاوتی هستند.
همچنین، حسگر ورودی و نقطهٔ پرس‌وجوی خروجی نقش یکسانی ندارند.
اصول \lr{DeepONet}، این جداسازی شاخه و تنه را روشن می‌کند.
\cite{lu2021deeponet}

اُنیشی و همکاران، جریان را تحت شرایط ورودی متفاوت بررسی کردند.
مطالعات تحول زمانی نیز، انتقال پاسخ میان بازه‌های زمانی را بررسی می‌کنند.
این دو مسیر، زمینهٔ مرتبطی برای عملگرهای این پژوهش فراهم می‌کنند.
اما فیزیک، خانوادهٔ ورودی و شیوهٔ ارزیابی آن‌ها باید حفظ شود.
\cite{Onishi2026FlowConditions,Wang2023LongTime}

در پژوهش حاضر، ورودی تابعی در دو مسئله نقش متفاوتی دارد.
در مسئلهٔ جریان، پروفایل ورودی به میدان سرعت و فشار مرتبط می‌شود.
حرکت ذره، سپس با استفاده از میدان سرعت محاسبه می‌شود.
در مسئلهٔ مغناطیسی، تحریک زمانی به پاسخ حرکتی متصل می‌شود.
هندسه و مقیاس‌های اختصاصی هر مسئله، در فصل سوم مشخص می‌شوند.
این دو تعمیم، به معنی پاسخ معتبر برای هر ورودی دلخواه نیستند.

پرسش اصلی، دامنهٔ اعتبار پاسخ در خانوادهٔ تعریف‌شده است.
برای میدان، مؤلفه‌های سرعت، فشار و باقیمانده باید جدا ارزیابی شوند.
برای پاسخ زمانی، خطای پنجره و خطای زنجیره هر دو اهمیت دارند.
دقت حرکت ذره نیز، جای ارزیابی کامل میدان را نمی‌گیرد.
پس استفادهٔ معکوس از عملگر، به اعتبارسنجی مستقیم آن وابسته می‌ماند.

\textbf{بازسازی جریان از مسیرهای محدود}

وان و همکاران، پیوند مسیر و میدان را با ردیاب‌ها بررسی کردند.
ژو و گراور، ویژگی‌های ذرات اینرسی‌دار را نیز در تخمین وارد کردند.
این مطالعات نشان می‌دهند نوع ذره و مشاهده، بخش اصلی صورت‌بندی است.
\cite{Wan2026TrajectoryFlowNet,ZhouGrauer2026Joint}

در مسئلهٔ حاضر، جریان از خانواده‌ای مشخص از ورودی‌ها انتخاب می‌شود.
هندسه و اطلاعات معلوم ذره، دامنهٔ مجهول‌ها را محدود می‌کنند.
عملگر مستقیم، طی تخمین ورودی ثابت نگه داشته می‌شود.
پس بازسازی با چند مسیر، در همین خانواده بررسی می‌شود.
این نتیجه، بازیابی هر میدان دلخواه از همان تعداد مسیر را اثبات نمی‌کند.

زمان‌داربودن مسیر، بخشی از اطلاعات مشاهده است.
در حالت بدون زمان، ترتیب و هندسه حفظ می‌شوند.
ساعت پنهان نیز همراه مجهول‌های جریان تخمین زده می‌شود.
اگر سرعت اولیه معلوم باشد، این اطلاع همچنان در مسئله باقی می‌ماند.
بنابراین، دو حالت باید با اطلاعات جانبی دقیق مقایسه شوند.
معیار هندسی مشترک، برای مقایسهٔ بازپخش آن‌ها لازم است.
خطای پروفایل، میدان و مسیر نیز باید از یکدیگر جدا بمانند.

پرسش این محور، کفایت مسیرها برای تعیین مجهول‌های خانوادهٔ منتخب است.
مکان مسیرها، مدت مشاهده و حساسیت پاسخ، در این کفایت نقش دارند.
تعداد کم مسیر، به‌تنهایی معیار برتری یا اثبات یکتایی نیست.
همچنین، خطای جانشین می‌تواند در تخمین ورودی جبران شود.
این مسئله، ضرورت ارزیابی با مرجع مستقل را تقویت می‌کند.
\cite{Cao2023ResidualCorrection}

\textbf{طراحی تحریک و هم‌ارزی فرمان‌ها}

مطالعات مغناطیسی، سطوح متفاوتی از ورودی را بررسی می‌کنند.
نیروی مؤثر، دامنهٔ منبع و فرمان سخت‌افزار، کمیت‌های یکسانی نیستند.
رهگیری در سطح نیرو و تخمین تجربی توزیع نیرو، نمونه‌های مرتبط هستند.
هیچ‌کدام، به‌تنهایی بازیابی دو فرمان زمانی را نتیجه نمی‌دهند.
\cite{ElHamdaoui2026Tracking,Strayer2024ForceMapping}

در صورت‌بندی دومنبعی حاضر، میدان از دو شکل فضایی ثابت ساخته می‌شود.
دامنه‌های زمانی، مجهول‌های تحریک را مشخص می‌کنند.
در مدل القایی، نیرو به گرادیان مربع بزرگی میدان وابسته است.
بنابراین، علامت مشترک دو فرمان می‌تواند از حرکت نامشخص بماند.
پس از فاصلهٔ مشترک صفر، ابهام زمانی بیشتری نیز ممکن است.
تنها اصلاح علامت سراسری، همهٔ هم‌ارزی‌های زمانی را حذف نمی‌کند.

پرسش این محور، طراحی تحریک مجاز برای بازتولید مسیر هدف است.
جریان معلوم، کالیبراسیون ذره و محدودیت قطبیت، اطلاعات مسئله را تغییر می‌دهند.
در طراحی، بازپخش مسیر هدف مبنای ارزیابی فرمان است.
همچنین، نتیجهٔ تک‌منبعی با فرض قطبیت معلوم، به دومنبعی تعمیم نمی‌یابد.
تحلیل حساسیت و نمایهٔ درست‌نمایی، ابزارهای مکمل بررسی محدودیت‌ها هستند.
بااین‌حال، نتیجهٔ عددی این ابزارها، باید همراه فرض‌هایشان تفسیر شود.
\cite{Raue2009ProfileLikelihood}

\textbf{پیوند پیشینه با صورت‌بندی و ارزیابی پژوهش}

جدول \ref{tab:lit-synthesis-thesis-links} ارتباط محورهای مرور و مدل‌های فصل سوم را خلاصه می‌کند.
این جدول، تفاوت صورت‌بندی را ثبت می‌کند و رتبه‌بندی دقت نیست.

\begin{table}[htbp]
\centering
\small
\setlength{\tabcolsep}{4pt}
\caption{پیوند محورهای مرور با صورت‌بندی پژوهش حاضر}
\label{tab:lit-synthesis-thesis-links}
\begin{tabular}{p{0.19\textwidth}p{0.30\textwidth}p{0.39\textwidth}}
\hline
محور & پیشینهٔ مرتبط & صورت‌بندی و آزمون متناظر \\
\hline
تشکیل آرایش
& هدایت بازخوردی و اپتوسیالی \cite{Shenoy2016StokesTrap,Erben2024Optofluidic}
& پیش‌بینی مسیر، طراحی رهاسازی و آزمون فاصله؛ بخش \mbox{\ref{sec:theory-particle-formation}} \\
عملگر جریان
& یادگیری جریان با ورودی متفاوت \cite{Onishi2026FlowConditions}
& پروفایل به میدان؛ ارزیابی مستقل؛ بخش \mbox{\ref{sec:forward-flow-operator}} \\
عملگر حرکت
& یادگیری تحول کوتاه‌مدت \cite{Wang2023LongTime}
& تحریک به مسیر؛ آزمون زنجیره؛ بخش \mbox{\ref{sec:forward-magnetic-operator}} \\
بازسازی جریان
& بازسازی لاگرانژی–اویلری \cite{Wan2026TrajectoryFlowNet,ZhouGrauer2026Joint}
& ورودی و ساعت پنهان؛ بخش‌های \mbox{\ref{sec:inverse-flow-timed}} و \mbox{\ref{sec:inverse-flow-geometry}} \\
طراحی معکوس تحریک
& رهگیری نیرو و تخمین توزیع نیرو \cite{ElHamdaoui2026Tracking,Strayer2024ForceMapping}
& مسیر هدف، تحریک مجاز و بازپخش مستقل؛ بخش \mbox{\ref{sec:inverse-magnetic-drive}} \\
\hline
\end{tabular}
\end{table}

جایگاه پژوهش، از ترکیب همین پرسش‌ها و محدودیت‌های مشخص تعریف می‌شود.
تفاوت صورت‌بندی، باید از برتری عددی متمایز بماند.
برتری عددی، به آزمون هم‌شرایط و معیارهای مشترک نیاز دارد.
کمبود نمونهٔ دقیقاً مشابه، برای ادعای تقدم کافی نیست.
در فصل سوم، مجهول‌ها، مدل‌ها و شیوهٔ آموزش مشخص می‌شوند.
فصل چهارم، شواهد پیش‌بینی، بازسازی و بازپخش را ارائه می‌کند.
تفسیر نهایی، باید به دامنهٔ همین شواهد و فرض‌ها محدود بماند.
